\documentclass[aps,prb,onecolumn,nofootinbib,superscriptaddress]{revtex4-2}

\usepackage{amsmath,amssymb,amsthm,mathtools,bm}
\usepackage{booktabs}
\usepackage[colorlinks=true,linkcolor=blue,citecolor=blue,urlcolor=blue]{hyperref}
\usepackage{microtype}

\newcommand{\I}{\mathbf{1}}
\newcommand{\E}{\mathbb{E}}
\newcommand{\Prob}{\mathbb{P}}
\newcommand{\Tr}{\operatorname{Tr}}
\newcommand{\diag}{\operatorname{diag}}
\newcommand{\rank}{\operatorname{rank}}
\newcommand{\spec}{\operatorname{spec}}
\newcommand{\artanh}{\operatorname{artanh}}
\newcommand{\dd}{\mathrm{d}}
\newcommand{\ket}[1]{\lvert #1\rangle}
\newcommand{\bra}[1]{\langle #1\rvert}
\newcommand{\cA}{\mathcal{A}}
\newcommand{\cC}{\mathcal{C}}
\newcommand{\cF}{\mathcal{F}}
\newcommand{\cH}{\mathcal{H}}
\newcommand{\cK}{\mathcal{K}}
\newcommand{\cM}{\mathcal{M}}
\newcommand{\cT}{\mathcal{T}}
\newcommand{\cX}{\mathcal{X}}
\newcommand{\bxi}{{\boldsymbol{\xi}}}

\theoremstyle{plain}
\newtheorem{proposition}{Proposition}

\theoremstyle{definition}
\newtheorem{definition}{Definition}
\newtheorem{remark}{Remark}

\makeatletter
\let\@bibdataout@init\relax
\let\@bibdataout@rev\relax
\makeatother

\begin{document}

\title{Free-Fermion Trajectory Lyapunov Spectroscopy with Gain and Loss}
\author{Asadullah Bhuiyan}
\date{\today}

\begin{abstract}
This note spells out how the trajectory-transfer construction used to extract
operator scaling dimensions at measurement-induced transitions simplifies for
number-conserving Gaussian fermion circuits.  The many-body Kraus operator of a
fixed measurement record is replaced, wherever possible, by one-body
determinants, covariance updates, and singular values of one-body products.
This simplification is not merely computational: it also exposes which
``transfer matrix'' is being diagonalized.  We distinguish Born-sampled
covariance trajectories, the trajectory-averaged mean channel, cumulative
Gaussian/Kraus operator data, Choi particle and hole transfer charts, and the
tangent cocycle used in finite-time Lyapunov simulations.  Particular care is
given to gain and loss.  Single-fermion creation and annihilation branches are
odd Gaussian amplitudes rather than ordinary invertible even Gaussian
operators; exact projected branches therefore generate zero or infinite
directions in one transfer chart.  The correct finite diagnostic is the
finite-sector particle or hole spectrum together with the null-sector
occupation and multiplicity.  With this bookkeeping in place, the usual
cylinder formulas for the effective central charge, scaling dimensions, and
multifractal moment exponents can be stated without conflating distinct
operators.
\end{abstract}

\maketitle
\setcounter{tocdepth}{2}
\tableofcontents

\section{Introduction}

Zabalo \emph{et al.} showed that a monitored circuit trajectory can be treated
as a sample of a disordered statistical mechanics model by the identification
\begin{equation}
  Z_\bxi = p_\bxi ,
  \label{eq:intro-Z-p}
\end{equation}
where \(\bxi\) is the measurement record and \(p_\bxi\) is its Born
probability~\cite{Zabalo2022}.  The logarithm of \(p_\bxi\) is then a sample
free energy, and the subleading singular values of the trajectory Kraus
operator play the role of transfer-matrix eigenvalues.  At a scale-invariant
point, finite-size corrections of these Lyapunov data give the effective
central charge and low-lying scaling dimensions of the corresponding
non-unitary conformal field theory (CFT).

For an interacting random circuit, this is a many-body calculation.  A
trajectory operator acts on a Hilbert space of dimension exponential in system
size, and its leading singular values must be extracted by a many-body
iteration.  In a number-conserving free-fermion circuit, the same conceptual
construction has a much sharper form.  A trajectory built from Gaussian
operations remains Gaussian, so its probabilities are determinants and its
finite transfer spectra are controlled by one-particle matrices or by
covariance-matrix Schur complements.  This is the sense in which the
Zabalo--Gullans--Huse--Pixley spectroscopy becomes a free-fermion Lyapunov
problem.

The purpose of this note is not to claim that every finite-time Lyapunov gap
currently measured in the class-A adaptive circuit is already a universal CFT
number.  The purpose is to make the route precise enough that a numerical
extraction cannot accidentally identify different objects.  This matters
especially in the adaptive circuit, where local feedback contains particle
gain and loss.  Those branches are perfectly legitimate Gaussian trajectory
amplitudes, but they live on singular boundaries of the usual invertible
particle-transfer chart.

The universal objects used below are standard but come from several
literatures.  Cylinder finite-size scaling and operator dimensions are reviewed
in Refs.~\cite{Cardy1986,DiFrancesco1997}; entanglement central charge is a
different use of conformal data~\cite{CalabreseCardy2004}.  Quenched
Lyapunov spectra are governed by multiplicative ergodic theory
\cite{Oseledets1968,FurstenbergKesten1960}, while multifractal spectra at
disordered critical points are reviewed in Ref.~\cite{EversMirlin2008}.  The
free-fermion monitored-circuit context is developed in
Refs.~\cite{JianBauerKeselmanLudwig2022,AlbertonBuchholdDiehl2021,Fava2023,
Poboiko2023,ChahineBuchhold2024}, and the topological Lyapunov/boundary-mode
language relevant to class A appears in
Refs.~\cite{LudwigSchulzBaldesStolz2013,XiaoKawabata2024,BhuiyanPanJian2026}.

\section{Object Dictionary}
\label{sec:dictionary}

The word ``transfer'' is overloaded.  Table~\ref{tab:dictionary} fixes the
objects used below.

\begin{table}[h]
\caption{\label{tab:dictionary}
Distinct objects that should not be identified without an explicit map.}
\begin{ruledtabular}
\begin{tabular}{p{0.23\linewidth}p{0.31\linewidth}p{0.34\linewidth}}
Object & Definition & What it can diagnose \\
\colrule
Born trajectory \(G_\bxi(t)\) &
State covariance evolved by the sampled record \(\bxi\) &
Trajectory-resolved nonlinear observables such as entropy, Chern markers, and
correlators \\
Mean channel \(\bar G(t)\) &
Deterministic covariance obtained after Born averaging the channel &
Linear observables of the average state and mixed-state spectral skeletons \\
Cumulative Kraus/Gaussian operator &
Unnormalized amplitude \(\hat K_\bxi(t)\) or its one-body representative when
the branch is even and regular &
Born probability, operator singular values, and trajectory free energy \\
Choi particle/hole transfer data &
Doubled Gaussian covariance of \(\hat K_\bxi(t)\), represented in particle or
hole charts plus endpoint data &
Finite-sector singular values and singular gain/loss null sectors \\
Tangent cocycle &
Linearization of \(G_{t+1}=F_{\alpha_t}(G_t)\) along a realized trajectory &
Stability exponents of covariance perturbations; near-marginal directions \\
CFT data &
Universal finite-size coefficients of the trajectory ensemble at a critical
point &
\(c_{\rm eff}\), scaling dimensions, and multifractal spectra \\
\end{tabular}
\end{ruledtabular}
\end{table}

The distinction between the third and fifth rows is essential.  The cumulative
Kraus or Choi transfer object asks for singular values of a realized operator.
The tangent cocycle asks how an infinitesimal state perturbation evolves along
the nonlinear stochastic covariance map.  They may display the same slow
physical sector, but they are not the same operator.

\section{Gaussian State Conventions}
\label{sec:gaussian-state-conventions}

Let \(c_i,c_i^\dagger\), \(i=1,\ldots,N\), be complex fermions on the
one-particle Hilbert space \(\cH\simeq \mathbb C^N\).  For a Gaussian state
\(\rho\), define
\begin{equation}
  C_{ij}=\Tr(\rho\,c_i^\dagger c_j),
  \qquad
  G=2C-\I .
  \label{eq:C-G-definition}
\end{equation}
For a pure Slater determinant with occupied one-particle frame
\(Q\in\mathbb C^{N\times q}\), \(Q^\dagger Q=\I_q\),
\begin{equation}
  \ket{\Psi_Q}=
  \prod_{a=1}^{q} c^\dagger(Q_a)\ket{0},
  \qquad
  C=QQ^\dagger .
  \label{eq:slater-frame}
\end{equation}
Hence
\begin{equation}
  C^2=C,
  \qquad
  G^2=(2C-\I)^2=\I .
  \label{eq:pure-identities}
\end{equation}
Conversely, a Hermitian \(C\) with \(C^2=C\) is the one-body projector onto
the occupied subspace of a fixed-charge pure Gaussian state.  The covariance
language is therefore equivalent to the Slater-frame language, but it is more
convenient for measurements and feedback.

\section{Occupation Measurements as Determinantal Updates}
\label{sec:occupation-measurements}

Let
\begin{equation}
  \ket{\chi}\in\cH,
  \qquad
  P=\ket{\chi}\bra{\chi},
  \qquad
  Q_\chi=\I-P,
  \qquad
  \hat n_\chi=\hat\chi^\dagger\hat\chi ,
  \label{eq:rank-one-mode}
\end{equation}
with \(\|\chi\|=1\).  The two projective measurement outcomes are
\(\hat n_\chi=1\) and \(\hat n_\chi=0\).  Their Born probabilities follow
immediately from Eq.~\eqref{eq:C-G-definition}:
\begin{equation}
  p_1=\Tr(PC)=\bra{\chi}C\ket{\chi},
  \qquad
  p_0=1-p_1=\Tr[P(\I-C)] .
  \label{eq:single-mode-born}
\end{equation}

The conditional covariance can be derived by choosing a basis in which
\(\chi\) is the first orbital and writing
\begin{equation}
  C=
  \begin{pmatrix}
    c & v^\dagger\\
    v & R
  \end{pmatrix},
  \qquad
  c=\bra{\chi}C\ket{\chi}.
  \label{eq:C-block-for-conditioning}
\end{equation}
For the occupied outcome, Wick contraction conditioned on \(n_\chi=1\) gives
\begin{equation}
  C^{(1)}_{\rm rest}=R-\frac{vv^\dagger}{c}.
  \label{eq:occupied-schur}
\end{equation}
For the empty outcome, the determinantal complement gives
\begin{equation}
  C^{(0)}_{\rm rest}=R+\frac{vv^\dagger}{1-c}.
  \label{eq:empty-schur}
\end{equation}
Restoring basis-independent notation,
\begin{align}
  C^{(1)}
  &=
  P+Q_\chi C Q_\chi
  -
  \frac{Q_\chi C P C Q_\chi}{\Tr(PC)},
  \label{eq:C-occupied-update}\\
  C^{(0)}
  &=
  Q_\chi C Q_\chi
  +
  \frac{Q_\chi C P C Q_\chi}{1-\Tr(PC)} .
  \label{eq:C-empty-update}
\end{align}
The cross blocks with the measured mode vanish because the measured mode is
now a classical definite occupation.

For several commuting measured modes with projector \(P_S\), the same
calculation gives determinant minors.  If \(S\) denotes the set of modes
conditioned to be occupied, then
\begin{equation}
  \Prob(S\ {\rm occupied})
  =
  \det(C_S),
  \label{eq:minor-occupied}
\end{equation}
where \(C_S=P_S C P_S\) restricted to \(\operatorname{Ran}P_S\).  Empty
conditions are obtained by replacing \(C\) with \(\I-C\), or by using the
usual determinantal point-process conditional kernel.  This is the first
free-fermion simplification of \(Z_\bxi=p_\bxi\): the trajectory weight is
assembled from determinant factors rather than from a many-body tensor
contraction.

\section{Even Gaussian Branches: Normalization and One-Body Singular Values}
\label{sec:even-gaussian-branches}

Consider an even number-conserving Gaussian branch whose action on creation
operators is represented by a one-body matrix \(A\):
\begin{equation}
  \hat K_A\,c^\dagger(\phi)\,\hat K_A^{-1}
  =
  c^\dagger(A\phi),
  \label{eq:even-branch-conjugation}
\end{equation}
with \(A\) invertible for the moment.  Acting on a Gaussian density matrix
with covariance \(C\), the unnormalized output \(\hat K_A\rho\hat K_A^\dagger\)
has normalization
\begin{equation}
  p_A(C)
  =
  \Tr(\hat K_A\rho\hat K_A^\dagger)
  =
  \det\!\left(\I-C+ACA^\dagger\right),
  \label{eq:gaussian-branch-determinant}
\end{equation}
up to the scalar prefactor used to normalize \(\hat K_A\).  The normalized
output covariance is
\begin{equation}
  C'
  =
  ACA^\dagger
  \left(\I-C+ACA^\dagger\right)^{-1}.
  \label{eq:gaussian-covariance-update}
\end{equation}
Thus the single-particle version of the many-body trace
\(\Tr(\hat K_A\rho\hat K_A^\dagger)\) is not a one-body trace.  It is the
determinant in Eq.~\eqref{eq:gaussian-branch-determinant}.  In the
\(G=2C-\I\) convention this same normalization is
\begin{equation}
  p_A(G)
  =
  \det\!\left[
    \frac{\I-G}{2}
    +
    A\frac{\I+G}{2}A^\dagger
  \right].
  \label{eq:gaussian-branch-determinant-G}
\end{equation}
With the convention of Eq.~\eqref{eq:even-branch-conjugation}, the mixed-state
formula is \(\det(\I-C+ACA^\dagger)\).  One should not replace it by
\(\det(\I-C+C A^\dagger A)\) in a mixed state unless an additional convention
or commutation relation justifies the reordering.  The \(A^\dagger A\) form
appears cleanly after restricting to the occupied Slater frame, as in
Eq.~\eqref{eq:pure-slater-even-branch} below.
Equation~\eqref{eq:gaussian-covariance-update} is fixed by two checks.  First,
if \(A\) is unitary, it reduces to \(C'=ACA^\dagger\).  Second, in a basis
where both \(C\) and \(A^\dagger A\) are diagonal, each occupation mode is
reweighted by the ordinary Bernoulli factor
\begin{equation}
  n_a'=\frac{s_a^2 n_a}{1-n_a+s_a^2 n_a},
  \label{eq:bernoulli-reweighting}
\end{equation}
where \(s_a\) is the corresponding singular value of \(A\).

For a pure Slater determinant, Eq.~\eqref{eq:slater-frame} gives the simpler
Gram-determinant formula
\begin{equation}
  p_A(Q)=\det\!\left(Q^\dagger A^\dagger A Q\right),
  \qquad
  C'=
  A Q\left(Q^\dagger A^\dagger A Q\right)^{-1}Q^\dagger A^\dagger ,
  \label{eq:pure-slater-even-branch}
\end{equation}
provided \(AQ\) has rank \(q\).  This formula is just the norm and projector
of the transported occupied frame.

For a trajectory \(\bxi=(\alpha_1,\ldots,\alpha_t)\) made only of regular even
branches, the one-body product is
\begin{equation}
  A_\bxi(t)=A_{\alpha_t}\cdots A_{\alpha_1}.
  \label{eq:one-body-product}
\end{equation}
If
\begin{equation}
  s_a[A_\bxi(t)]\sim \exp\!\left(t\gamma_a\right),
  \label{eq:one-body-lyapunov}
\end{equation}
then the many-body Fock-space singular values of the second-quantized operator
are products of the one-body singular values.  Indeed, singular-value
decomposing \(A_\bxi=U\diag(s_1,\ldots,s_N)V^\dagger\), the second-quantized
unitaries \(\Gamma(U)\) and \(\Gamma(V)\) do not change singular values, while
\(\Gamma(\diag s)\) acts on a Fock basis state \(\ket{I}\) as
\begin{equation}
  \Gamma(\diag s)\ket{I}
  =
  \left(\prod_{a\in I}s_a\right)\ket{I}.
  \label{eq:fock-product-singular}
\end{equation}
Thus
\begin{equation}
  \lambda_I
  =
  \lim_{t\to\infty}\frac{1}{t}
  \log \sigma_I^2
  =
  2\sum_{a\in I}\gamma_a ,
  \label{eq:many-body-additive-lyapunov}
\end{equation}
apart from any trajectory-independent scalar normalization convention.  In a
fixed-charge sector, the allowed sets \(I\) have fixed cardinality.

\subsection{One-Body Spectra Versus Fock-Space Spectra}
\label{sec:one-body-many-body-spectra}

It is useful to state the spectral dictionary without CFT language.  Let
\(A_\bxi(t)\) be any finite one-body Gaussian transfer matrix associated with
a fixed trajectory and a fixed particle/hole chart.  Write its singular values
as
\begin{equation}
  s_1(t)\ge s_2(t)\ge\cdots\ge s_N(t)>0,
  \qquad
  \kappa_a=\lim_{t\to\infty}\frac{1}{t}\log s_a(t).
  \label{eq:one-body-kappa}
\end{equation}
The induced operator on the \(q\)-particle sector is the exterior power
\(\bigwedge^q A_\bxi(t)\).  Its singular values are
\begin{equation}
  s_I^{(q)}(t)
  =
  \prod_{a\in I}s_a(t),
  \qquad
  I\subset\{1,\ldots,N\},
  \qquad |I|=q,
  \label{eq:exterior-power-singular-values}
\end{equation}
and its Lyapunov exponents are
\begin{equation}
  \kappa_I^{(q)}
  =
  \sum_{a\in I}\kappa_a .
  \label{eq:exterior-power-exponents}
\end{equation}
Thus the many-body spectrum is not a new diagonalization problem once the
one-body singular data are known, but neither is it identical to the one-body
list.  It is the additive spectrum of the appropriate exterior-power or
Fock-space representation.

The universal gaps are gaps relative to the dominant many-body configuration
in the sector being measured.  If \(I_0\) maximizes \(\kappa_I^{(q)}\) in the
chosen \(q\)-particle sector, then
\begin{equation}
  \mu_{I\leftarrow I_0}^{(q)}
  =
  \kappa_{I_0}^{(q)}-\kappa_I^{(q)}
  =
  \sum_{a\in I_0}\kappa_a-\sum_{a\in I}\kappa_a
  \ge 0 .
  \label{eq:many-body-gap-difference-of-sums}
\end{equation}
For example, a single particle-hole replacement \(a\in I_0\to b\notin I_0\)
has
\begin{equation}
  \mu_{a\to b}^{(q)}=\kappa_a-\kappa_b .
  \label{eq:particle-hole-gap}
\end{equation}
Only in a convention where the low-energy one-body data have already been
converted to positive excitation rapidities
\(\epsilon_\nu\) measured relative to \(I_0\) does the many-body tower take the
simple form
\begin{equation}
  \mu_A=\sum_\nu n_\nu(A)\epsilon_\nu .
  \label{eq:excitation-rapidity-tower}
\end{equation}
This is the convention used by the finite-gap estimator in the CPU analysis
script.  The integer pattern \(n_\nu(A)\) is not fixed by linear algebra alone:
it is fixed by the charge, parity, particle/hole chart, boundary condition,
and observable sector.  The default choice ``add the lowest one positive
rapidity'' is therefore a sector assumption, not a theorem.

If the trajectory transfer carries a scalar normalization
\(c_\bxi(t)\Gamma(A_\bxi(t))\), then \(t^{-1}\log|c_\bxi(t)|\) shifts every
many-body exponent in the sector by the same amount.  It affects absolute
free-energy densities but cancels from the gaps
\(\mu_{I\leftarrow I_0}^{(q)}\).  This is why CFT dimensions are extracted
from many-body gaps, while \(c_{\rm eff}\) is extracted separately from the
trajectory probability free energy.

\section{Gain and Loss: Odd Branches and Singular Transfer Charts}
\label{sec:gain-loss}

The adaptive class-A circuit also uses feedback that supplies or removes a
fermion from a measured mode.  These operations are not ordinary even
invertible Gaussian branches.  On a pure Slater state they are
\begin{equation}
  \hat K_+=\hat\chi^\dagger,
  \qquad
  \hat K_-=\hat\chi ,
  \label{eq:odd-state-branches}
\end{equation}
possibly composed with surrounding even Gaussian evolution.  If the branch
has nonzero amplitude, creation and annihilation update the physical covariance
by
\begin{align}
  C_{\rm add}
  &=
  C+
  \frac{(\I-C)P(\I-C)}
       {\Tr[(\I-C)P]},
  \label{eq:addition-update}\\
  C_{\rm rem}
  &=
  C-
  \frac{CPC}
       {\Tr(CP)} .
  \label{eq:removal-update}
\end{align}
Equation~\eqref{eq:addition-update} says that creation fills the component of
\(\chi\) that was empty; Eq.~\eqref{eq:removal-update} says that annihilation
removes the component of \(\chi\) that was occupied.  The denominators are the
corresponding branch probabilities.
\begin{equation}
  p_+(C)=\Tr[P(\I-C)]=1-\bra{\chi}C\ket{\chi},
  \qquad
  p_-(C)=\Tr(PC)=\bra{\chi}C\ket{\chi}.
  \label{eq:gain-loss-probabilities}
\end{equation}
In the \(G=2C-\I\) convention these become
\begin{equation}
  p_+(G)=\frac{1-\bra{\chi}G\ket{\chi}}{2},
  \qquad
  p_-(G)=\frac{1+\bra{\chi}G\ket{\chi}}{2}.
  \label{eq:gain-loss-probabilities-G}
\end{equation}

\subsection{How Odd Branches Compose}
\label{sec:odd-composition}

The important point for transfer matrices is that gain and loss do compose
with the transfer product, but not as ordinary multiplication by another
invertible one-body matrix in the same charge sector.  Let
\(\cH_1\) be the one-particle Hilbert space and let \(\Gamma(A)\) denote the
second-quantized even Gaussian operator induced by a one-body map \(A\).  On
the \(q\)-particle sector \(\bigwedge^q\cH_1\), define exterior multiplication
and contraction by
\begin{align}
  \varepsilon_\chi \Psi &= \chi\wedge\Psi,
  &
  \varepsilon_\chi&:\bigwedge^q\cH_1\to\bigwedge^{q+1}\cH_1,
  \label{eq:exterior-creation-map}\\
  \iota_\chi(v_1\wedge\cdots\wedge v_q)
  &=
  \sum_{a=1}^{q}(-1)^{a-1}
  \langle\chi|v_a\rangle\,
  v_1\wedge\cdots\widehat{v_a}\cdots\wedge v_q,
  &
  \iota_\chi&:\bigwedge^q\cH_1\to\bigwedge^{q-1}\cH_1 .
  \label{eq:interior-annihilation-map}
\end{align}
Here \(\widehat v_a\) means that \(v_a\) is omitted.  The odd Kraus branches
are precisely these charge-changing maps, up to the scalar normalization and
fermion-ordering convention:
\begin{equation}
  \hat\chi^\dagger \leftrightarrow \varepsilon_\chi,
  \qquad
  \hat\chi \leftrightarrow \iota_\chi .
  \label{eq:odd-branch-exterior-interior}
\end{equation}
They obey the exact composition identities
\begin{align}
  \Gamma(A)\,\varepsilon_\chi
  &=
  \varepsilon_{A\chi}\,\Gamma(A),
  \label{eq:even-creation-composition}\\
  \iota_\chi\,\Gamma(A)
  &=
  \Gamma(A)\,\iota_{A^\dagger\chi}.
  \label{eq:even-annihilation-composition}
\end{align}
Thus, for two even segments \(A_1,A_2\), a gain insertion between them gives
\begin{equation}
  \Gamma(A_2)\,\varepsilon_\chi\,\Gamma(A_1)
  =
  \varepsilon_{A_2\chi}\,\Gamma(A_2A_1),
  \label{eq:gain-insertion-composition}
\end{equation}
while a loss insertion gives
\begin{equation}
  \Gamma(A_2)\,\iota_\chi\,\Gamma(A_1)
  =
  \Gamma(A_2A_1)\,\iota_{A_1^\dagger\chi}.
  \label{eq:loss-insertion-composition}
\end{equation}
Equations~\eqref{eq:gain-insertion-composition} and
\eqref{eq:loss-insertion-composition} are the mechanical reason gain/loss
cannot be represented as a harmless extra factor in the same one-body transfer
matrix product.  The trajectory transfer is a word in even maps and
charge-changing exterior/interior maps,
\begin{equation}
  \widehat{\cT}_\bxi
  =
  \Gamma(A_m)\,\mathcal O_m\,
  \Gamma(A_{m-1})\,\mathcal O_{m-1}\cdots
  \mathcal O_1\,\Gamma(A_0),
  \qquad
  \mathcal O_j\in\{\varepsilon_{\chi_j},\iota_{\chi_j}\},
  \label{eq:odd-even-transfer-word}
\end{equation}
so its many-body singular spectrum is a charge-sector-changing Fock-space
spectrum.  If one insists on a fixed particle-transfer chart, each
\(\varepsilon_\chi\) or \(\iota_\chi\) forces a rank-one direction to an
endpoint or to a null direction, and the remaining finite block is the
reduced transfer matrix used below.  The missing endpoint occupation is not
lost information; it is the record of which exterior/interior map occurred.

\subsection{Several Odd Insertions: Same-Type Blocks and Mixed Words}
\label{sec:several-odd-insertions}

The single-insertion discussion is enough to identify the obstruction, but it
is not enough for an actual adaptive trajectory.  A trajectory can contain
several gains, several losses, or a mixed ordered string of both.  The
correct rule is not to commute the odd events into a convenient order.  The
correct rule is to keep the chronological exterior-algebra word.

First consider a block containing only gains.  For
\begin{equation}
  \widehat{\cT}_{+}
  =
  \Gamma(A_r)\varepsilon_{\chi_r}
  \Gamma(A_{r-1})\varepsilon_{\chi_{r-1}}\cdots
  \varepsilon_{\chi_1}\Gamma(A_0),
  \label{eq:many-gain-word}
\end{equation}
successive use of Eq.~\eqref{eq:even-creation-composition} gives
\begin{equation}
  \widehat{\cT}_{+}
  =
  \varepsilon_{\eta_r}\varepsilon_{\eta_{r-1}}\cdots
  \varepsilon_{\eta_1}\,
  \Gamma(A_{\rm tot}),
  \qquad
  A_{\rm tot}=A_rA_{r-1}\cdots A_0,
  \label{eq:many-gain-normal-form}
\end{equation}
with transported gain modes
\begin{equation}
  \eta_j=A_rA_{r-1}\cdots A_j\chi_j .
  \label{eq:transported-gain-modes}
\end{equation}
Thus the \(q\)-particle action is
\begin{equation}
  \Psi\in\bigwedge^q\cH_1
  \quad\longmapsto\quad
  \eta_r\wedge\eta_{r-1}\wedge\cdots\wedge\eta_1
  \wedge A_{\rm tot}^{\wedge q}\Psi
  \in\bigwedge^{q+r}\cH_1 .
  \label{eq:many-gain-sector-map}
\end{equation}
Permuting two gain insertions changes the sign of the exterior product, and
with interleaved even maps it also changes the transported vectors.  Therefore
only a block of adjacent gains in the same chart can be regarded as
order-independent up to an overall sign.

For adjacent gains acting on a state with covariance \(C\), put the measured
modes in the matrix
\begin{equation}
  X=(\chi_1,\ldots,\chi_r),
  \qquad
  E_X=X^\dagger(\I-C)X .
  \label{eq:gain-block-empty-gram}
\end{equation}
Wick's theorem gives the Born probability
\begin{align}
  p_{+X}(C)
  &=
  \Tr\!\left[
  \hat\chi_r^\dagger\cdots\hat\chi_1^\dagger
  \rho
  \hat\chi_1\cdots\hat\chi_r
  \right]
  =
  \det E_X .
  \label{eq:many-gain-born}
\end{align}
When \(E_X\) is nonsingular, the conditional covariance is
\begin{equation}
  C_{+X}
  =
  C+(\I-C)X\,E_X^{-1}X^\dagger(\I-C).
  \label{eq:many-gain-covariance}
\end{equation}
This is the multi-mode version of Eq.~\eqref{eq:addition-update}: the branch
fills the \(r\)-dimensional component of the requested gain subspace that
lies in the empty subspace of the incoming Gaussian state.  If
\(\det E_X=0\), the branch is Pauli blocked in at least one linear
combination; the exact trajectory weight is zero and the corresponding
transfer chart is rank deficient.

The loss-only block is dual.  For
\begin{equation}
  \widehat{\cT}_{-}
  =
  \Gamma(A_r)\iota_{\chi_r}
  \Gamma(A_{r-1})\iota_{\chi_{r-1}}\cdots
  \iota_{\chi_1}\Gamma(A_0),
  \label{eq:many-loss-word}
\end{equation}
successive use of Eq.~\eqref{eq:even-annihilation-composition} gives
\begin{equation}
  \widehat{\cT}_{-}
  =
  \Gamma(A_{\rm tot})\,
  \iota_{\zeta_r}\iota_{\zeta_{r-1}}\cdots\iota_{\zeta_1},
  \label{eq:many-loss-normal-form}
\end{equation}
where
\begin{equation}
  \zeta_j=(A_{j-1}A_{j-2}\cdots A_0)^\dagger\chi_j .
  \label{eq:transported-loss-modes}
\end{equation}
For adjacent losses in the current state chart,
\begin{equation}
  O_X=X^\dagger C X,
  \label{eq:loss-block-occupied-gram}
\end{equation}
and
\begin{align}
  p_{-X}(C)
  &=
  \Tr\!\left[
  \hat\chi_r\cdots\hat\chi_1
  \rho
  \hat\chi_1^\dagger\cdots\hat\chi_r^\dagger
  \right]
  =
  \det O_X,
  \label{eq:many-loss-born}\\
  C_{-X}
  &=
  C-CX\,O_X^{-1}X^\dagger C .
  \label{eq:many-loss-covariance}
\end{align}
The determinant is now the occupied-subspace Gram determinant.  A zero
determinant means that at least one requested removal direction has no
occupied component left after the previous removals.

The genuinely delicate case is a mixed string.  Let
\begin{equation}
  d_j=
  \begin{cases}
    \hat\chi_j^\dagger, & \sigma_j=+,\\
    \hat\chi_j, & \sigma_j=-,
  \end{cases}
  \qquad
  \hat K_{\bm\sigma}=d_r d_{r-1}\cdots d_1 .
  \label{eq:mixed-odd-word}
\end{equation}
The exterior and interior maps satisfy the Clifford relations
\begin{equation}
  \{\varepsilon_u,\varepsilon_v\}=0,
  \qquad
  \{\iota_u,\iota_v\}=0,
  \qquad
  \iota_u\varepsilon_v+\varepsilon_v\iota_u=\langle u|v\rangle .
  \label{eq:exterior-interior-clifford-relations}
\end{equation}
The last relation is the whole non-commuting issue.  Moving a loss past a gain
does not merely generate a sign; it generates a scalar contraction.  In the
simplest same-mode example,
\begin{equation}
  \hat\chi\hat\chi^\dagger=1-\hat n_\chi,
  \qquad
  \hat\chi^\dagger\hat\chi=\hat n_\chi .
  \label{eq:same-mode-gain-loss-projectors}
\end{equation}
A gain followed by a loss projects onto the initially empty sector, whereas a
loss followed by a gain projects onto the initially occupied sector.  These
two branches have probabilities \(1-n_\chi\) and \(n_\chi\), respectively;
they are not two equivalent endpoint insertions.

For an arbitrary mixed ordered string the compact formula is a Wick/Pfaffian
formula.  Define the ordered list of linear fermion operators appearing in
\(\hat K_{\bm\sigma}^\dagger\hat K_{\bm\sigma}\),
\begin{equation}
  L_1,\ldots,L_{2r}
  =
  d_1^\dagger,\ldots,d_r^\dagger,d_r,\ldots,d_1 .
  \label{eq:mixed-kdagger-k-list}
\end{equation}
For an even gauge-invariant Gaussian state, the only nonzero elementary
contractions are
\begin{equation}
  \langle \hat c^\dagger(u)\hat c(v)\rangle
  =
  \langle u|C|v\rangle,
  \qquad
  \langle \hat c(u)\hat c^\dagger(v)\rangle
  =
  \langle u|(\I-C)|v\rangle .
  \label{eq:number-conserving-wick-contractions}
\end{equation}
Let \(W\) be the \(2r\times2r\) antisymmetric matrix
\begin{equation}
  W_{\mu\nu}=
  \begin{cases}
    \langle L_\mu L_\nu\rangle, & \mu<\nu,\\
    -\langle L_\nu L_\mu\rangle, & \mu>\nu,\\
    0, & \mu=\nu .
  \end{cases}
  \label{eq:mixed-wick-matrix}
\end{equation}
Then
\begin{equation}
  p_{\bm\sigma}(C)
  =
  \Tr(\hat K_{\bm\sigma}\rho\hat K_{\bm\sigma}^\dagger)
  =
  \langle L_1\cdots L_{2r}\rangle
  =
  \operatorname{Pf} W .
  \label{eq:mixed-word-born-pfaffian}
\end{equation}
Equations~\eqref{eq:mixed-kdagger-k-list}--\eqref{eq:mixed-word-born-pfaffian}
are the mechanical prescription for non-commuting gain/loss words.  The
chronological order appears in the ordered list \(L_\mu\); changing the order
changes \(W\), and hence changes both the Born weight and the conditional
Gaussian state.

The post-branch covariance is obtained from the same Wick expansion.  In a
Majorana basis, let \(\gamma_a\) be the physical Majoranas and let
\(\mathcal G_{ab}\) denote the antisymmetrized Gaussian two-point contraction.
With
\begin{equation}
  B_{\mu a}=\langle L_\mu\gamma_a\rangle,
  \label{eq:mixed-schur-B}
\end{equation}
the conditional antisymmetrized contraction is the Pfaffian Schur complement
\begin{equation}
  \mathcal G'_{ab}
  =
  \mathcal G_{ab}
  +
  \sum_{\mu,\nu=1}^{2r}
  B_{\mu a}(W^{-1})_{\mu\nu}B_{\nu b},
  \qquad
  \operatorname{Pf}W\ne0 .
  \label{eq:mixed-pfaffian-schur-update}
\end{equation}
Equivalently, one can evaluate the sequential rank-one updates
Eqs.~\eqref{eq:addition-update} and \eqref{eq:removal-update} in the actual
time order.  The Pfaffian formula is useful because it makes explicit that
the answer is an ordered Gaussian conditioning problem, not a product of
commuting scalar probabilities.

For transfer spectroscopy the conclusion is sharp.  Same-type blocks may be
compressed into determinant updates and a rank-\(r\) endpoint/null sector,
provided the corresponding Gram matrix has rank \(r\).  A mixed non-commuting
block must either be kept as an ordered exterior/interior Clifford word or
embedded in a regularized ordered Majorana/Pin product.  In either
description, the finite singular values are read only after this ordered
reduction; endpoint multiplicities are ranks of the resulting constrained
subspaces, not simply the number of gain/loss events.

\subsection{\texorpdfstring{Transfer-Matrix Group: \(O(2N,\mathbb C)\) and
the \(U(1)\) Levi Chart}{Transfer-Matrix Group: O(2N,C) and the U(1) Levi Chart}}
\label{sec:nambu-transfer-group-structure}

Before introducing a regulator, it is useful to state the group constraint.
Let
\begin{equation}
  \Psi=\begin{pmatrix}c\\ c^\dagger\end{pmatrix},
  \qquad
  \Omega=\begin{pmatrix}0&\I\\ \I&0\end{pmatrix},
  \qquad
  \{\Psi_\mu,\Psi_\nu\}=\Omega_{\mu\nu}.
  \label{eq:nambu-group-car-form}
\end{equation}
If an invertible Gaussian operator \(\hat V\) acts by
\begin{equation}
  \hat V\Psi\hat V^{-1}=\mathbb T[\hat V]\Psi ,
  \label{eq:nambu-transfer-action-group}
\end{equation}
then the canonical anticommutation relations force
\begin{equation}
  \mathbb T[\hat V]\,\Omega\,\mathbb T[\hat V]^T=\Omega .
  \label{eq:nambu-transfer-orthogonality}
\end{equation}
Thus the unprotected Nambu transfer matrix lives in the complex orthogonal
group \(O(2N,\mathbb C)\), not in a free \(\mathrm{GL}(2N,\mathbb C)\).  In a
Majorana basis this is the ordinary complex orthogonal constraint; the complex
Nambu basis only rewrites the same bilinear form as \(\Omega\).

For a number-conserving invertible even branch the Nambu transfer is block
diagonal,
\begin{equation}
  \mathbb T=
  \begin{pmatrix}
    T_h&0\\
    0&T_p
  \end{pmatrix},
  \qquad
  T_hT_p^T=\I,
  \qquad
  T_p=T_h^{-T}.
  \label{eq:u1-levi-nambu-transfer}
\end{equation}
The \(U(1)\)-equivariant invertible transfer group is therefore the
\(\mathrm{GL}(N,\mathbb C)\) Levi chart parametrized by \(T_p\).  The
particle/hole inverse-transpose relation is not an additional symmetry
assumption; it is Eq.~\eqref{eq:nambu-transfer-orthogonality} restricted to
the \(U(1)\)-commuting block form.

Exact gain and loss sit outside this invertible \(U(1)\) chart for a simple
reason.  A charge-homogeneous linear insertion is a null vector of the CAR
form:
\begin{equation}
  \hat\chi^\dagger=w_\chi^T\Psi,
  \quad
  w_\chi=\begin{pmatrix}0\\ \chi\end{pmatrix},
  \qquad
  \hat\chi=\bar u_\chi^T\Psi,
  \quad
  \bar u_\chi=\begin{pmatrix}\chi^*\\0\end{pmatrix},
  \qquad
  w_\chi^T\Omega w_\chi=\bar u_\chi^T\Omega\bar u_\chi=0 .
  \label{eq:gain-loss-null-covectors}
\end{equation}
Null linear Clifford elements are nilpotent and cannot be inverted.  Hence no
invertible regulator for a single gain or loss can remain \(U(1)\)-equivariant:
the regulator must mix the two charge sectors, run in \(O(2N,\mathbb C)\), and
return to the charge-homogeneous operator only as a singular endpoint.


\subsection{\texorpdfstring{Moving Nambu Charts and the \(O(2N,\mathbb C)\)
Regularization}{Moving Nambu Charts and the O(2N,C) Regularization}}
\label{sec:onambu-regularization}

There is another useful way to say the same thing.  The fixed particle chart is
singular because creation and annihilation are nilpotent maps.  If instead one
works in the full Nambu or Majorana Clifford algebra, an odd branch can be
approximated by an invertible Clifford element.  For a normalized mode
\(\chi\), define
\begin{equation}
  v_\varepsilon
  =
  \sqrt{\varepsilon}\,\hat\chi
  +
  \frac{1}{\sqrt{\varepsilon}}\,\hat\chi^\dagger,
  \qquad
  \varepsilon>0 .
  \label{eq:pin-regularized-odd-element}
\end{equation}
Using \(\hat\chi^2=(\hat\chi^\dagger)^2=0\) and
\(\{\hat\chi,\hat\chi^\dagger\}=1\), one finds
\begin{equation}
  v_\varepsilon^2=1,
  \qquad
  v_\varepsilon^{-1}=v_\varepsilon .
  \label{eq:pin-element-invertible}
\end{equation}
Thus \(v_\varepsilon\) is not a creation operator.  It is an odd Clifford, or
Pin, element whose conjugation action on Nambu operators is an
\(O(2N,\mathbb C)\) transformation.  On the distinguished mode,
\begin{equation}
  v_\varepsilon\,\hat\chi\,v_\varepsilon^{-1}
  =
  \frac{1}{\varepsilon}\hat\chi^\dagger,
  \qquad
  v_\varepsilon\,\hat\chi^\dagger\,v_\varepsilon^{-1}
  =
  \varepsilon \hat\chi ,
  \label{eq:pin-conjugation-single-mode}
\end{equation}
while modes anticommuting with \(\chi\) are transformed by the corresponding
orthogonal reflection, with signs fixed by the chosen Clifford ordering.
Equation~\eqref{eq:pin-conjugation-single-mode} shows exactly what the fixed
particle chart calls a divergence: in the full Nambu chart it is a
particle-hole reflection accompanied by a diagonal gauge
\(\varepsilon^{\pm1}\).

Pure gain and loss are recovered only as projective singular limits.  For
gain,
\begin{equation}
  \sqrt{\varepsilon}\,v_\varepsilon
  =
  \varepsilon\hat\chi+\hat\chi^\dagger
  \xrightarrow{\varepsilon\to0}
  \hat\chi^\dagger ,
  \label{eq:pin-gain-limit}
\end{equation}
whereas for loss one may equivalently use
\begin{equation}
  \frac{1}{\sqrt{\varepsilon}}\,v_\varepsilon
  =
  \hat\chi+\frac{1}{\varepsilon}\hat\chi^\dagger
  \label{eq:pin-loss-bad-scaling}
\end{equation}
with \(\varepsilon\to\infty\), or exchange
\(\varepsilon\leftrightarrow\varepsilon^{-1}\) in
Eq.~\eqref{eq:pin-regularized-odd-element}.  More symmetrically, define
\begin{equation}
  w_\varepsilon
  =
  \frac{1}{\sqrt{\varepsilon}}\,\hat\chi
  +
  \sqrt{\varepsilon}\,\hat\chi^\dagger ,
  \qquad
  w_\varepsilon^2=1,
  \qquad
  \sqrt{\varepsilon}\,w_\varepsilon
  \xrightarrow{\varepsilon\to0}
  \hat\chi .
  \label{eq:pin-loss-limit}
\end{equation}
The singular gain/loss maps are therefore boundary points of an invertible
Pin/Bogoliubov cocycle, rather than failures of the underlying Clifford
description.

This regularization is useful but not free.  The physical trajectory
probability is not the projective \(O(2N,\mathbb C)\) action alone.  For the
unrescaled operator \(v_\varepsilon\),
\begin{align}
  \Tr(v_\varepsilon\rho v_\varepsilon^\dagger)
  &=
  \varepsilon\,\Tr(\hat\chi\rho\hat\chi^\dagger)
  +
  \varepsilon^{-1}\Tr(\hat\chi^\dagger\rho\hat\chi)
  \nonumber\\
  &=
  \varepsilon\,p_-(C)+\varepsilon^{-1}p_+(C),
  \label{eq:pin-regularized-branch-weight}
\end{align}
because the same-mode nilpotent cross terms contain \(\hat\chi^2\) or \((\hat\chi^\dagger)^2\).
After multiplying by the projective scalar \(\sqrt{\varepsilon}\) appropriate
to the gain limit, Eq.~\eqref{eq:pin-regularized-branch-weight} tends to
\(p_+(C)\).  After using the loss scaling in
Eq.~\eqref{eq:pin-loss-limit}, it tends to \(p_-(C)\).  Consequently the
moving \(O(2N,\mathbb C)\) chart regularizes the transfer direction, but the
scalar cocycle must be kept separately to recover the Born weight and the
trajectory free energy.

In practical terms, the two descriptions are equivalent bookkeeping choices:
\begin{equation}
  \begin{array}{c}
  \hbox{fixed particle/hole chart}
  \end{array}
  \quad\Longleftrightarrow\quad
  \begin{array}{c}
  \hbox{finite block plus endpoint/null sectors}
  \end{array},
  \label{eq:fixed-chart-bookkeeping}
\end{equation}
whereas
\begin{equation}
  \begin{array}{c}
  O(2N,\mathbb C)\hbox{ moving Nambu chart}
  \end{array}
  \quad\Longleftrightarrow\quad
  \begin{array}{c}
  \hbox{finite Pin/Bogoliubov cocycle plus divergent gauge}\\
  \hbox{and scalar normalization}
  \end{array}.
  \label{eq:moving-chart-bookkeeping}
\end{equation}
The moving chart may be the cleaner representation for composing long
trajectories with gain and loss, especially if one wants a single nonsingular
Nambu transfer product.  It does not eliminate the gain/loss complication; it
moves it from endpoint eigenvalues into the gauge/scalar part of the cocycle.

\subsection{Projective Measurements, Gauge Renormalization, and Particle/Hole
Reciprocity}
\label{sec:projector-gauge-reciprocity}

The same singular-chart issue already appears for ordinary projective
occupation measurements.  These are even branches, so they do not mix particle
and hole sectors at finite regulator.  Let \(P=\chi\chi^\dagger\) be the
one-body projector onto the measured particle wavefunction and define
\begin{equation}
  D_\varepsilon(P)=\I-P+\varepsilon P .
  \label{eq:projector-regulator-D}
\end{equation}
The occupied and empty projectors can be approached by the invertible even
branches
\begin{equation}
  \hat \Pi_{1,\varepsilon}
  =
  \hat n_\chi+\varepsilon(1-\hat n_\chi),
  \qquad
  \hat \Pi_{0,\varepsilon}
  =
  (1-\hat n_\chi)+\varepsilon\hat n_\chi .
  \label{eq:regularized-occupation-projectors}
\end{equation}
For \(\varepsilon>0\), these are honest Gaussian transfer matrices and
therefore compose by ordinary matrix multiplication with every other
finite-regulator Gaussian factor.  In the Nambu convention
\(\Psi=(c,c^\dagger)^T\), their block-diagonal transfer matrices are
\begin{align}
  \mathbb T_{1,\varepsilon}
  &=
  \begin{pmatrix}
    D_\varepsilon(P)^T & 0\\
    0 & D_\varepsilon(P)^{-1}
  \end{pmatrix},
  \label{eq:occupied-projector-nambu-transfer}\\
  \mathbb T_{0,\varepsilon}
  &=
  \begin{pmatrix}
    D_\varepsilon(P)^{-T} & 0\\
    0 & D_\varepsilon(P)
  \end{pmatrix}.
  \label{eq:empty-projector-nambu-transfer}
\end{align}
Equivalently, on the measured mode,
\begin{align}
  \hat \Pi_{1,\varepsilon}\hat\chi
  \hat \Pi_{1,\varepsilon}^{-1}
  &=
  \varepsilon\hat\chi,
  &
  \hat \Pi_{1,\varepsilon}\hat\chi^\dagger
  \hat \Pi_{1,\varepsilon}^{-1}
  &=
  \varepsilon^{-1}\hat\chi^\dagger,
  \label{eq:occupied-projector-mode-action}\\
  \hat \Pi_{0,\varepsilon}\hat\chi
  \hat \Pi_{0,\varepsilon}^{-1}
  &=
  \varepsilon^{-1}\hat\chi,
  &
  \hat \Pi_{0,\varepsilon}\hat\chi^\dagger
  \hat \Pi_{0,\varepsilon}^{-1}
  &=
  \varepsilon\hat\chi^\dagger .
  \label{eq:empty-projector-mode-action}
\end{align}
Thus an occupied measurement produces a particle-sector pole and a hole-sector
zero on the measured mode, while an empty measurement produces the opposite
zero/pole pair.  At every nonzero \(\varepsilon\) the particle and hole blocks
are still related by CAR preservation:
\begin{equation}
  A_h(\varepsilon)A_p(\varepsilon)^T=\I,
  \qquad
  A_h(\varepsilon)=A_p(\varepsilon)^{-T}.
  \label{eq:finite-epsilon-particle-hole-inverse}
\end{equation}
The exact projector limit destroys the full-rank chart, not the finite
\(\varepsilon\) relation.  On the active complement of \(P\), the reduced
particle and hole blocks remain inverse-transpose partners; on the measured
direction one instead records an endpoint occupation.

The same regularized projector is equally transparent in Majorana language.
Define
\(\gamma_{1,\chi}=\hat\chi+\hat\chi^\dagger\) and
\(\gamma_{2,\chi}=-i(\hat\chi-\hat\chi^\dagger)\).  Then
\begin{equation}
  \hat n_\chi=\frac{1+i\gamma_{1,\chi}\gamma_{2,\chi}}{2}.
  \label{eq:number-projector-majoranas}
\end{equation}
Let \(B_\chi=i\gamma_{1,\chi}\gamma_{2,\chi}\), so \(B_\chi^2=1\), and put
\begin{equation}
  \eta_\varepsilon=\frac{1}{2}\log\frac{1}{\varepsilon}.
  \label{eq:projector-boost-eta}
\end{equation}
Then
\begin{equation}
  \hat\Pi_{1,\varepsilon}
  =
  \sqrt{\varepsilon}\,e^{\eta_\varepsilon B_\chi},
  \qquad
  \hat\Pi_{0,\varepsilon}
  =
  \sqrt{\varepsilon}\,e^{-\eta_\varepsilon B_\chi}.
  \label{eq:projector-majorana-boost}
\end{equation}
The scalar \(\sqrt{\varepsilon}\) is part of the Kraus normalization and drops
out of conjugation.  The occupied branch acts on
\((\gamma_{1,\chi},\gamma_{2,\chi})\) by the complex orthogonal boost
\begin{equation}
  R_{1,\varepsilon}^{(\chi)}
  =
  \begin{pmatrix}
    \cosh(2\eta_\varepsilon)&-i\sinh(2\eta_\varepsilon)\\
    i\sinh(2\eta_\varepsilon)&\cosh(2\eta_\varepsilon)
  \end{pmatrix},
  \label{eq:occupied-projector-majorana-boost}
\end{equation}
and the empty branch uses \(R_{0,\varepsilon}^{(\chi)}
=(R_{1,\varepsilon}^{(\chi)})^{-1}\).  The Nambu operators
\(\hat\chi\) and \(\hat\chi^\dagger\) are the light-cone eigenvectors of this
Majorana boost, with eigenvalues \(\varepsilon\) and
\(\varepsilon^{-1}\) for the occupied outcome and the reciprocal eigenvalues
for the empty outcome.  Thus the Majorana and Nambu derivations are the same
singular transfer written in different bases: Majorana sees a divergent boost,
while Nambu diagonalizes that boost into particle/hole zero-pole factors.

This also explains what can and cannot be achieved by ``renormalizing'' the
transfer matrix.  Multiplying the Kraus operator by a scalar changes the Born
weight but cancels out of the conjugation action.  Multiplying the whole
Nambu matrix by a scalar is not a legitimate CAR-preserving transfer matrix,
except for the signs \(r=\pm1\).  A legitimate number-conserving Nambu gauge
has the reciprocal form
\begin{equation}
  \mathbb G(R)=
  \begin{pmatrix}
    R^{-T}&0\\
    0&R
  \end{pmatrix},
  \qquad
  \mathbb G(R)\Omega\mathbb G(R)^T=\Omega .
  \label{eq:nambu-reciprocal-gauge}
\end{equation}
For example,
\begin{equation}
  \mathbb T_{1,\varepsilon}=\mathbb G[D_\varepsilon(P)^{-1}],
  \qquad
  \mathbb T_{0,\varepsilon}=\mathbb G[D_\varepsilon(P)] .
  \label{eq:projectors-as-singular-gauges}
\end{equation}
One may use such gauges to center a numerical calculation or to subtract a
known regulator slope, but this does not produce a unique finite full
transfer matrix in the exact limit.  It only moves the zero/pole into the
chart transition.  The invariant data are the finite active block, the
removed regulator exponents, and the endpoint occupation.

Gain and loss are the odd analogue.  Their regulated Nambu matrices in
Eqs.~\eqref{eq:nambu-regulated-gain-matrix} and
\eqref{eq:nambu-regulated-loss-matrix} are not block diagonal: the zero/pole
pair appears in the particle-hole mixing blocks instead of in the diagonal
projector blocks.  Nevertheless the finite-\(\varepsilon\) statement is the
same.  The full Nambu or Majorana transfer is an ordinary
\(O(2N,\mathbb C)\) matrix product.  The exact \(\varepsilon\to0\) limit is a
singular chart limit whose finite content is a reduced particle/hole transfer
matrix plus a filled or empty endpoint sector.

The resulting interpretation of the whole monitored Gaussian circuit is:
choose a finite regulator for every exact projector, gain, and loss branch;
compose the associated Nambu or Majorana one-body matrices by literal ordered
matrix multiplication; keep the scalar determinant/Pfaffian cocycle for the
Born weight; and only then take the regulator limit.  The limit returns the
exact many-body Gaussian branch, but the single-particle representative is a
singular-limit object rather than a globally full-rank particle-transfer
matrix.

\subsection{Majorana Transfer Form of Exact and Regularized Gain/Loss}
\label{sec:majorana-gain-loss-transfer}

The preceding discussion can be made completely explicit in a Majorana basis.
For the mode \(\chi\), define two Majorana operators
\begin{equation}
  \gamma_{1,\chi}=\hat\chi+\hat\chi^\dagger,
  \qquad
  \gamma_{2,\chi}=-i(\hat\chi-\hat\chi^\dagger),
  \label{eq:mode-majoranas}
\end{equation}
so that
\begin{equation}
  \hat\chi=\frac{\gamma_{1,\chi}+i\gamma_{2,\chi}}{2},
  \qquad
  \hat\chi^\dagger=\frac{\gamma_{1,\chi}-i\gamma_{2,\chi}}{2}.
  \label{eq:creation-annihilation-majoranas}
\end{equation}
Thus exact loss and exact gain are linear Clifford insertions,
\begin{equation}
  \hat K_-=\hat\chi=a_-\cdot\gamma,
  \qquad
  \hat K_+=\hat\chi^\dagger=a_+\cdot\gamma,
  \label{eq:exact-gain-loss-majorana-linear}
\end{equation}
with, in the two-dimensional subspace spanned by
\((\gamma_{1,\chi},\gamma_{2,\chi})\),
\begin{equation}
  a_-=\frac{1}{2}(1,i),
  \qquad
  a_+=\frac{1}{2}(1,-i).
  \label{eq:null-majorana-vectors}
\end{equation}
These vectors are null with respect to the complex orthogonal form:
\begin{equation}
  a_\pm\cdot a_\pm=0,
  \qquad
  \hat K_\pm^2=0 .
  \label{eq:exact-gain-loss-null}
\end{equation}
This is the exact Majorana-basis obstruction.  Gain and loss are
Majorana-linear operators, but because they are nilpotent they do not define
an invertible conjugation action
\(\gamma_i\mapsto \hat K_\pm\gamma_i\hat K_\pm^{-1}\).  Exact gain/loss is
therefore a null Clifford insertion, not an element of
\(O(2N,\mathbb C)\).

The regularized branch
\begin{equation}
  v_\varepsilon
  =
  \sqrt{\varepsilon}\,\hat\chi
  +
  \frac{1}{\sqrt{\varepsilon}}\,\hat\chi^\dagger
  =
  a_\varepsilon\cdot\gamma
  \label{eq:majorana-regularized-branch}
\end{equation}
has
\begin{equation}
  a_\varepsilon
  =
  \frac{1}{2}
  \left(
  \sqrt{\varepsilon}+\varepsilon^{-1/2},
  i\sqrt{\varepsilon}-i\varepsilon^{-1/2}
  \right),
  \qquad
  a_\varepsilon\cdot a_\varepsilon=1 .
  \label{eq:majorana-regularized-vector}
\end{equation}
Hence \(v_\varepsilon^2=1\), as in
Eq.~\eqref{eq:pin-element-invertible}, and \(v_\varepsilon\) induces an
honest complex-orthogonal Majorana transfer matrix
\begin{equation}
  v_\varepsilon\,\gamma_i\,v_\varepsilon^{-1}
  =
  \sum_j R_{ij}^{(\varepsilon)}\gamma_j,
  \qquad
  R^{(\varepsilon)}\in O(2N,\mathbb C).
  \label{eq:regularized-majorana-transfer}
\end{equation}
Thus the precise statement is
\begin{align}
  \hbox{exact gain/loss}
  &=
  \hbox{null Majorana-linear Clifford insertion},
  \nonumber\\
  \hbox{regularized gain/loss}
  &=
  \hbox{invertible Pin element inducing }O(2N,\mathbb C)\hbox{ transfer}.
  \label{eq:exact-versus-regularized-majorana}
\end{align}

The Born probabilities make clear what is, and is not, contained in this
regularized transfer matrix.  Let
\begin{equation}
  n_\chi=\bra{\chi}C\ket{\chi}
  =
  \frac{1+\bra{\chi}G\ket{\chi}}{2}.
  \label{eq:n-chi-C-G}
\end{equation}
For exact loss and gain,
\begin{equation}
  p_-(C)
  =
  \Tr(\hat\chi\rho\hat\chi^\dagger)
  =
  n_\chi,
  \qquad
  p_+(C)
  =
  \Tr(\hat\chi^\dagger\rho\hat\chi)
  =
  1-n_\chi .
  \label{eq:exact-gain-loss-born}
\end{equation}
For a same-mode regularized odd branch
\begin{equation}
  \hat K=a\,\hat\chi+b\,\hat\chi^\dagger ,
  \label{eq:general-odd-same-mode-branch}
\end{equation}
an even Gaussian density matrix has no parity-odd expectation values, and
\(\hat\chi^2=(\hat\chi^\dagger)^2=0\).  The cross terms therefore vanish and
\begin{equation}
  p_K(C)
  =
  \Tr(\hat K\rho \hat K^\dagger)
  =
  |a|^2 n_\chi+|b|^2(1-n_\chi).
  \label{eq:regularized-odd-born-general}
\end{equation}
For \(v_\varepsilon\) this gives
\begin{equation}
  p_{v_\varepsilon}(C)
  =
  \varepsilon n_\chi+\varepsilon^{-1}(1-n_\chi).
  \label{eq:regularized-v-born}
\end{equation}
This quantity diverges in the gain limit because \(v_\varepsilon\) contains a
divergent scalar gauge.  The projectively rescaled gain approximation
\begin{equation}
  \hat K_\varepsilon^+
  =
  \sqrt{\varepsilon}\,v_\varepsilon
  =
  \varepsilon\hat\chi+\hat\chi^\dagger
  \label{eq:regularized-gain-K}
\end{equation}
has
\begin{equation}
  p_\varepsilon^+(C)
  =
  \varepsilon^2 n_\chi+(1-n_\chi)
  \xrightarrow{\varepsilon\to0}
  1-n_\chi .
  \label{eq:regularized-gain-born-limit}
\end{equation}
Similarly, using
\begin{equation}
  \hat K_\varepsilon^-
  =
  \hat\chi+\varepsilon\hat\chi^\dagger
  \label{eq:regularized-loss-K}
\end{equation}
for loss gives
\begin{equation}
  p_\varepsilon^-(C)
  =
  n_\chi+\varepsilon^2(1-n_\chi)
  \xrightarrow{\varepsilon\to0}
  n_\chi .
  \label{eq:regularized-loss-born-limit}
\end{equation}
Therefore the \(O(2N,\mathbb C)\) transfer matrix regularizes the projective
Majorana dynamics, but the physical trajectory weight lives in the scalar
normalization used to take the singular gain/loss limit.  Dropping that scalar
would lose precisely the contribution needed for the trajectory free energy
and hence for \(c_{\rm eff}\).


\subsection{Regulated Products in the Majorana/Nambu Chart}
\label{sec:regulated-majorana-products}

The regulated Majorana chart simplifies composition at finite regulator, but
only after one separates the projective transfer direction from the scalar
normalization.  Let
\begin{equation}
  \gamma=(\gamma_1,\ldots,\gamma_{2N})^T,
  \qquad
  \{\gamma_a,\gamma_b\}=2\delta_{ab} .
  \label{eq:majorana-vector-convention}
\end{equation}
A regulated odd branch may be written as
\begin{equation}
  \hat K_j(\varepsilon)
  =
  c_j(\varepsilon)\,\hat u_j(\varepsilon),
  \qquad
  \hat u_j(\varepsilon)=u_j(\varepsilon)^T\gamma,
  \qquad
  u_j(\varepsilon)^T u_j(\varepsilon)=1 .
  \label{eq:regulated-branch-pin-factorization}
\end{equation}
The transpose in \(u^Tu\) is the complex orthogonal bilinear form, not the
Hermitian norm.  Since \(\hat u_j^2=1\), conjugation by \(\hat u_j\) is the
complex orthogonal reflection
\begin{equation}
  \hat u_j\gamma_a\hat u_j^{-1}
  =
  \sum_b R[u_j]_{ab}\gamma_b,
  \qquad
  R[u_j]_{ab}=-\delta_{ab}+2u_{j,a}u_{j,b},
  \qquad
  R[u_j]\in O(2N,\mathbb C).
  \label{eq:pin-reflection-matrix}
\end{equation}
Even Gaussian pieces also induce complex-orthogonal Majorana matrices.  Hence
a regulated trajectory with arbitrary non-commuting gains and losses becomes
an ordinary ordered product of even rotations and odd reflections.  The precise
ordering convention matters, because conjugation is an antihomomorphism; the
Nambu statement in Eqs.~\eqref{eq:nambu-antihomomorphism} and
\eqref{eq:arbitrary-regulated-nambu-product} fixes this convention explicitly.
This is the real simplification: there is no separate exterior/interior algebra
to commute by hand.  The price is that the resulting matrix is only the
projective transfer matrix.

For a single mode, the Nambu pair
\(\Psi_\chi=(\hat\chi,\hat\chi^\dagger)^T\) makes the regulator singularity
explicit.  The projectively rescaled gain and loss approximants
\begin{equation}
  \hat K^+_\varepsilon=\hat\chi^\dagger+\varepsilon\hat\chi,
  \qquad
  \hat K^-_\varepsilon=\hat\chi+\varepsilon\hat\chi^\dagger
  \label{eq:regulated-gain-loss-approximants}
\end{equation}
obey \((\hat K^\pm_\varepsilon)^2=\varepsilon\), so they are invertible for
\(\varepsilon\ne0\).  Their conjugation actions are
\begin{align}
  \hat K^+_\varepsilon\Psi_\chi(\hat K^+_\varepsilon)^{-1}
  &=
  \begin{pmatrix}
    0 & \varepsilon^{-1}\\
    \varepsilon & 0
  \end{pmatrix}
  \Psi_\chi,
  \label{eq:nambu-gain-transfer}\\
  \hat K^-_\varepsilon\Psi_\chi(\hat K^-_\varepsilon)^{-1}
  &=
  \begin{pmatrix}
    0 & \varepsilon\\
    \varepsilon^{-1} & 0
  \end{pmatrix}
  \Psi_\chi .
  \label{eq:nambu-loss-transfer}
\end{align}
Thus finite \(\varepsilon\) replaces a zero or pole in a fixed particle chart
by a particle-hole exchange with singular gauge factors
\(\varepsilon^{\pm1}\).  These factors are not finite critical exponents.

It is useful to write the arbitrary-word formula directly in the Nambu basis.
Let
\begin{equation}
  \Psi=\begin{pmatrix}c\\ c^\dagger\end{pmatrix},
  \qquad
  \Omega=\begin{pmatrix}0&\I\\ \I&0\end{pmatrix},
  \qquad
  \{\Psi_\mu,\Psi_\nu\}=\Omega_{\mu\nu} .
  \label{eq:nambu-car-form}
\end{equation}
For a linear form
\begin{equation}
  \hat\ell(w)=w^T\Psi,
  \qquad
  \hat\ell(w)^2=\frac{1}{2}w^T\Omega w,
  \label{eq:nambu-linear-form}
\end{equation}
the exact gain/loss operators are precisely null vectors,
\(w^T\Omega w=0\), and so have no inverse.  A finite regulator replaces that
null vector by a non-null vector.  If \(w^T\Omega w\ne0\), conjugation by the
normalized Pin element induces
\begin{equation}
  \hat\ell(w)\Psi\hat\ell(w)^{-1}
  =
  \mathbb R[w]\Psi,
  \qquad
  \mathbb R[w]
  =
  -\I+\frac{2\Omega w w^T}{w^T\Omega w},
  \qquad
  \mathbb R[w]\Omega\mathbb R[w]^T=\Omega .
  \label{eq:nambu-pin-reflection}
\end{equation}
The factor \(\Omega w\), rather than \(w\), appears because \(w\) labels a
linear functional \(w^T\Psi\).  Equation~\eqref{eq:nambu-pin-reflection} is
the Nambu version of the Majorana reflection in
Eq.~\eqref{eq:pin-reflection-matrix}.

For a mode \(\chi\), with
\(\hat\chi^\dagger=\sum_i\chi_i c_i^\dagger\) and
\(\hat\chi=\sum_i\chi_i^*c_i\), define
\begin{equation}
  w_+(\varepsilon)=
  \begin{pmatrix}\varepsilon^{1/2}\chi^*\\ \varepsilon^{-1/2}\chi\end{pmatrix},
  \qquad
  w_-(\varepsilon)=
  \begin{pmatrix}\varepsilon^{-1/2}\chi^*\\ \varepsilon^{1/2}\chi\end{pmatrix} .
  \label{eq:nambu-regulated-gain-loss-vectors}
\end{equation}
Then \(\hat\ell[w_+(\varepsilon)]\) tends projectively to gain as
\(\varepsilon\to0\), while \(\hat\ell[w_-(\varepsilon)]\) tends projectively
to loss.  Since \(w_\pm^T\Omega w_\pm=2\), the Nambu reflection matrices are
\begin{align}
  \mathbb R_+(\chi,\varepsilon)
  &=
  \begin{pmatrix}
    -(\I-P) & \varepsilon^{-1}\chi\chi^T\\
    \varepsilon\chi^*\chi^\dagger & -(\I-P^T)
  \end{pmatrix},
  \label{eq:nambu-regulated-gain-matrix}\\
  \mathbb R_-(\chi,\varepsilon)
  &=
  \begin{pmatrix}
    -(\I-P) & \varepsilon\chi\chi^T\\
    \varepsilon^{-1}\chi^*\chi^\dagger & -(\I-P^T)
  \end{pmatrix},
  \label{eq:nambu-regulated-loss-matrix}
\end{align}
where \(P=\chi\chi^\dagger\).  For \(N=1\) these reduce to
Eqs.~\eqref{eq:nambu-gain-transfer} and \eqref{eq:nambu-loss-transfer}.
The off-diagonal blocks are the charge-breaking regulator pieces.  The full
matrix remains exactly CAR-orthogonal; if one drops either the divergent or
vanishing off-diagonal block before taking the limit, the missing product
\(\varepsilon^{-1}\varepsilon P\) is precisely the occupation constraint that
has to be restored by hand in the fixed particle/hole chart.

This statement becomes concrete when the neighboring even segment is
number-conserving.  Write
\begin{equation}
  \mathbb A=
  \begin{pmatrix}
    A_h&0\\
    0&A_p
  \end{pmatrix},
  \qquad
  A_hA_p^T=\I,
  \label{eq:nambu-number-conserving-even-block}
\end{equation}
and set \(P_h=P=\chi\chi^\dagger\), \(P_p=P^T\).  For an odd insertion on
the right side of this even transfer, ordinary Nambu multiplication gives
\begin{align}
  \mathbb A\mathbb R_+(\chi,\varepsilon)
  &=
  \begin{pmatrix}
    -A_h(\I-P_h) & \varepsilon^{-1}A_h\chi\chi^T\\
    \varepsilon A_p\chi^*\chi^\dagger & -A_p(\I-P_p)
  \end{pmatrix},
  \label{eq:even-times-regulated-gain-blocks}\\
  \mathbb A\mathbb R_-(\chi,\varepsilon)
  &=
  \begin{pmatrix}
    -A_h(\I-P_h) & \varepsilon A_h\chi\chi^T\\
    \varepsilon^{-1}A_p\chi^*\chi^\dagger & -A_p(\I-P_p)
  \end{pmatrix}.
  \label{eq:even-times-regulated-loss-blocks}
\end{align}
Thus the diagonal particle/hole blocks have a finite, rank-deficient limit,
while one charge-mixing block vanishes and the other diverges.  For a right
insertion the fixed-chart finite blocks are
\begin{equation}
  A_h\mapsto A_h(\I-P_h),
  \qquad
  A_p\mapsto A_p(\I-P_p),
  \label{eq:right-insertion-fixed-chart-limit}
\end{equation}
up to the unimportant reflection sign.  For a left insertion the same
calculation gives left multiplication by the complementary projectors,
\begin{equation}
  A_h\mapsto(\I-P_h)A_h,
  \qquad
  A_p\mapsto(\I-P_p)A_p .
  \label{eq:left-insertion-fixed-chart-limit}
\end{equation}
Inside a longer word the mode appearing in \(P_h\) or \(P_p\) is transported
by the adjacent even matrices, exactly as in
Eqs.~\eqref{eq:gain-insertion-composition} and
\eqref{eq:loss-insertion-composition}.  The exact gain/loss operation is
therefore recovered from the regulated product as a reduced finite
particle/hole transfer plus endpoint occupation data: gain and loss have the
same finite rank drop in the diagonal blocks, but gain has the
\(\varepsilon^{-1}\) pole in the particle-hole block and records a filled
endpoint, while loss has the opposite pole and records an empty endpoint.  A
fixed particle chart or fixed hole chart sees this as a zero or infinity; the
full Nambu chart sees it as a singular gauge direction of an otherwise
ordinary \(O(2N,\mathbb C)\) matrix product.

Now let the exact trajectory word be an arbitrary Gaussian word made of even
Gaussian pieces and exact gain/loss insertions.  Its regulated version can be
written as
\begin{equation}
  \hat K_{\bxi,\varepsilon}
  =
  s_{\bxi}(\varepsilon)
  \hat E_m\hat\ell[w_m(\varepsilon)]\hat E_{m-1}\cdots
  \hat\ell[w_1(\varepsilon)]\hat E_0,
  \label{eq:arbitrary-regulated-gaussian-word}
\end{equation}
where each \(\hat E_j\) is an invertible even Gaussian operator, not assumed to
be number conserving, and \(s_{\bxi}(\varepsilon)\) is the scalar projective
normalization.  If
\begin{equation}
  \hat E_j\Psi\hat E_j^{-1}=\mathbb A_j\Psi,
  \qquad
  \mathbb A_j\Omega\mathbb A_j^T=\Omega,
  \label{eq:even-nambu-transfer}
\end{equation}
then, with the convention \(\hat V\Psi\hat V^{-1}=\mathbb T[\hat V]\Psi\),
the map \(\hat V\mapsto\mathbb T[\hat V]\) is an antihomomorphism:
\begin{equation}
  \mathbb T[\hat V_2\hat V_1]
  =
  \mathbb T[\hat V_1]\mathbb T[\hat V_2] .
  \label{eq:nambu-antihomomorphism}
\end{equation}
Therefore the regulated word has the transfer matrix
\begin{equation}
  \mathbb T_{\bxi,\varepsilon}
  =
  \mathbb A_0\mathbb R[w_1(\varepsilon)]\mathbb A_1\cdots
  \mathbb R[w_m(\varepsilon)]\mathbb A_m,
  \qquad
  \mathbb T_{\bxi,\varepsilon}\Omega\mathbb T_{\bxi,\varepsilon}^T=\Omega .
  \label{eq:arbitrary-regulated-nambu-product}
\end{equation}
This is the clean Nambu statement for non-commuting gain/loss strings: all
ordering is ordinary matrix ordering in Eq.~\eqref{eq:arbitrary-regulated-nambu-product}.
No separate rule allows the gain factors to be commuted past the loss factors.

\subsection{Laurent Limit of an Arbitrary Regulated Word}
\label{sec:regulated-word-laurent-limit}

The preceding product formula is exact but it is not yet the exact
gain/loss endpoint.  For every \(\varepsilon>0\) the regulated Nambu, or
Majorana, representative is an honest \(O(2N,\mathbb C)\) matrix and products
are literal ordered matrix products.  The singular endpoint is recovered only
after passing to a graph/covariance chart and separating the scalar
normalization.  Taking the entrywise \(\varepsilon\to0\) limit of a raw
particle or hole block is generally the wrong operation.

A useful way to organize an arbitrary exact word is by its running charge.
Read the odd insertions in chronological order and define
\begin{equation}
  Q_j=\#\{\hbox{gains among }1,\ldots,j\}
      -\#\{\hbox{losses among }1,\ldots,j\},
  \qquad
  Q_0=0 .
  \label{eq:running-charge-definition}
\end{equation}
Cut the word whenever \(Q_j\) returns to its running minimum.  These cuts
isolate neutral blocks that can close internally from charge-changing tails.
For a neutral block \(B\) with nonzero contraction amplitude, the exact
operator has the form
\begin{equation}
  \hat K_B
  =
  \lambda_B\,\hat K_{T_B},
  \qquad
  \lambda_B\ne0 ,
  \label{eq:balanced-block-normal-form}
\end{equation}
where \(\hat K_{T_B}\) is an even Gaussian operator with a finite one-body
transfer \(T_B\) on the active graph chart.  For a block with net charge
\(q>0\),
\begin{equation}
  \hat K_B
  =
  \lambda_B\,
  \hat c^\dagger(\xi_q)\cdots\hat c^\dagger(\xi_1)
  \hat K_{T_B},
  \label{eq:unbalanced-gain-block-normal-form}
\end{equation}
and for \(q<0\) the exterior product is replaced by an ordered interior
product of \(|q|\) annihilation modes.  Thus an unbalanced block is not a
matrix alone; it is endpoint orbitals plus one matrix plus a scalar.

For a balanced mixed block, the scalar and the finite graph transfer are
fixed by the causal loss--gain Gram.  Let \(i\) label gains with incoming
wavefunctions \(\alpha_i\), let \(j\) label losses with outgoing
wavefunctions \(\beta_j\), and let \(M_{ji}\) be the even one-body
propagation from gain \(i\) to loss \(j\) through the pieces lying between
them.  The internal contraction matrix is
\begin{equation}
  \widetilde G_{ji}
  =
  \Theta(t_j^\ell>t_i^g)\,
  \langle \beta_j|M_{ji}|\alpha_i\rangle .
  \label{eq:causal-loss-gain-gram}
\end{equation}
The causal step function is the algebraic memory of the operator order:
same-type pairings vanish by isotropy of the gain and loss null subspaces,
while a loss can only contract a gain that occurs earlier in the word.  If
\(\det\widetilde G\ne0\), then
\begin{equation}
  \lambda_B=\pm\det\widetilde G ,
  \label{eq:balanced-block-amplitude}
\end{equation}
where the sign is the fermionic ordering sign.  This scalar is part of the
Born-weight cocycle.  It is not determined by the projective one-body
transfer matrix.

After partitioning the finite-\(\varepsilon\) Nambu graph into external
propagating legs \(b\), future-gain legs \(fg\), and loss legs \(l\), the
Schur complement gives the dressed master formula
\begin{equation}
  \frac{T_B}{\lambda_B}
  =
  M_{bb}-W_{fg}\,\widetilde G^{-1}\,W_l^\dagger .
  \label{eq:balanced-block-master-formula}
\end{equation}
Equivalently, if the regulated one-body graph is written in block form with
the lower-right graph block \(A_{dd}(\varepsilon)\), the finite object that
converges is
\begin{equation}
  \lim_{\varepsilon\to0}
  \bigl(A_{dd}(\varepsilon)^T\bigr)^{-1}
  =
  M_{bb}-W_{fg}\,\widetilde G^{-1}\,W_l^\dagger .
  \label{eq:graph-block-convergence}
\end{equation}
This is the same dressed-projector formula obtained by direct Choi/covariance
contraction.  The raw particle block \(A_{cc}(\varepsilon)\) can diverge, and
its projectively rescaled limit is typically only a flag.  The block in
Eq.~\eqref{eq:graph-block-convergence}, not the raw particle block, is the
finite transfer chart relevant for active-sector singular data.

The same-mode identities are a useful check:
\begin{equation}
  \hat\chi\,\hat\chi^\dagger=\I-\hat n_\chi,
  \qquad
  \hat\chi^\dagger\hat\chi=\hat n_\chi .
  \label{eq:gain-loss-same-mode-closures-chart}
\end{equation}
They are the one-mode version of Eq.~\eqref{eq:balanced-block-master-formula}.
Same-type gain or loss blocks instead grow endpoint occupation constraints.
For gains the Born factor is the empty-space Gram
\(\det[X^\dagger(\I-C)X]\); for losses it is the occupied-space Gram
\(\det[Y^\dagger C Y]\).  Non-orthogonal wavefunctions are handled by these
Gram determinants and the corresponding Schur complements; small or zero
determinants signal redundant, blocked, or ill-conditioned exact branches,
not additional finite Lyapunov gaps.

The finite spectrum is extracted from singular-value asymptotics only after
this separation.  If \(s_a(\varepsilon)\) are singular values of the
regulated product, define
\begin{equation}
  \log s_a(\varepsilon)
  =
  \kappa_a\log\frac{1}{\varepsilon}
  +
  \rho_a
  +
  o(1).
  \label{eq:regulated-singular-value-asymptotics}
\end{equation}
CAR preservation pairs nonzero slopes \(\kappa_a\) with opposite slopes.
Those directions are endpoint or gauge sectors.  The regulator-stable
offsets \(\rho_a\) with \(\kappa_a=0\), together with the scalar
\(\lambda_B\) and the filled/empty endpoint orbitals, are the complete
singular data.  This is why the regulated Majorana basis simplifies
composition but does not remove endpoint or Born-weight bookkeeping.

\subsection{Case-by-Case Regulated Endpoint Limits}
\label{sec:regulated-endpoint-case-derivations}

Here is the same regulator limit written mechanically for the elementary
cases.  In this subsection the regulator is denoted by \(\delta\), to avoid
confusing it with the exterior multiplication operator
\(\varepsilon_x\).  Work first in a chart where all even pieces have already
transported the odd wavefunctions to the relevant edge.  The even pieces are
restored by replacing the displayed modes by the transported modes of
Eqs.~\eqref{eq:transported-gain-modes} and
\eqref{eq:transported-loss-modes}, and by multiplying the active block by the
total even one-body transfer.  Define
\begin{equation}
  \hat k_\delta^+(x)=\hat c^\dagger(x)+\delta\,\hat c(x),
  \qquad
  \hat k_\delta^-(y)=\hat c(y)+\delta\,\hat c^\dagger(y).
  \label{eq:projectively-normalized-gain-loss}
\end{equation}
For normalized \(x\) and \(y\), \((\hat k_\delta^\pm)^2=\delta\), so
\(\delta^{-1/2}\hat k_\delta^\pm\) is an invertible Pin element at every
\(\delta\ne0\).  Thus the finite-\(\delta\) one-body Nambu/Majorana
representative composes by ordinary matrix multiplication.  The exact branch
is the leading coefficient of the operator expansion as \(\delta\to0\).

\paragraph*{1. One gain.}
The regulated operator is
\begin{equation}
  \hat k_\delta^+(x)
  =
  \hat c^\dagger(x)+\delta\hat c(x)
  =
  \hat c^\dagger(x)+O(\delta).
  \label{eq:one-gain-delta-limit}
\end{equation}
Thus the exact limit is the exterior map
\(\varepsilon_x:\Psi_q\mapsto x\wedge\Psi_q\).  It is not a scalar and not a
same-charge transfer.  For an incoming gauge-invariant Gaussian covariance
\(C\),
\begin{equation}
  p_{+,x}(C)
  =
  \Tr\!\left[\hat c^\dagger(x)\rho\hat c(x)\right]
  =
  x^\dagger(\I-C)x .
  \label{eq:one-gain-delta-born}
\end{equation}
If \(x=\|x\|\hat x\), the exact operator is
\(\|x\|\hat c^\dagger(\hat x)\).  The scalar \(\|x\|\) belongs to the branch
amplitude, while the active particle chart removes the filled endpoint,
\begin{equation}
  T_{\rm act}^{+x}
  =
  T_{\rm even}(\I-P_{\hat x}),
  \qquad
  P_{\hat x}=\hat x\hat x^\dagger .
  \label{eq:one-gain-active-transfer}
\end{equation}
The finite-\(\delta\) \(O(2N,\mathbb C)\) matrix has one
\(\delta^{-1}\) direction and one reciprocal \(\delta\) direction in the
particle/hole pair of \(x\).  Those directions are the filled endpoint
sector, not finite transfer gaps.

\paragraph*{2. Two gains.}
For two gains in chronological order \(x_1\) then \(x_2\),
\begin{align}
  \hat k_\delta^+(x_2)\hat k_\delta^+(x_1)
  &=
  \hat c^\dagger(x_2)\hat c^\dagger(x_1)
  +\delta\!\left[
    \hat c^\dagger(x_2)\hat c(x_1)
    +\hat c(x_2)\hat c^\dagger(x_1)
  \right]
  +\delta^2\hat c(x_2)\hat c(x_1)
  \nonumber\\
  &=
  \hat c^\dagger(x_2)\hat c^\dagger(x_1)
  +O(\delta).
  \label{eq:two-gain-regulated-expansion}
\end{align}
The middle term contains the contraction
\(\hat c(x_2)\hat c^\dagger(x_1)=
\langle x_2|x_1\rangle-\hat c^\dagger(x_1)\hat c(x_2)\), but it is still
suppressed by \(\delta\).  Therefore, when
\(x_1\wedge x_2\ne0\), the exact endpoint is a two-dimensional filled
subspace.  The Born factor is
\begin{equation}
  p_{++,X}(C)
  =
  \det
  \begin{pmatrix}
    x_1^\dagger(\I-C)x_1 & x_1^\dagger(\I-C)x_2\\
    x_2^\dagger(\I-C)x_1 & x_2^\dagger(\I-C)x_2
  \end{pmatrix},
  \qquad
  X=(x_1,x_2).
  \label{eq:two-gain-born-determinant}
\end{equation}
If \(X=UR\) is a QR decomposition of the independent endpoint span, then
\begin{equation}
  \hat c^\dagger(x_2)\hat c^\dagger(x_1)
  =
  \det(R)\,
  \hat c^\dagger(u_2)\hat c^\dagger(u_1),
  \qquad
  P_X=UU^\dagger .
  \label{eq:two-gain-qr-limit}
\end{equation}
The active transfer is \(T_{\rm even}(\I-P_X)\), and the Nambu singular
spectrum contains two reciprocal endpoint pairs.  If
\(\operatorname{rank}X<2\), the exact two-gain branch vanishes by
antisymmetry; the regulator then selects a \(\delta\)-suppressed different
word and the exact branch should be treated as blocked or censored.

\paragraph*{3. \(N\) gains.}
For \(N\) gains, put \(X=(x_1,\ldots,x_N)\) and write
\begin{equation}
  \hat K_{+,N}(\delta)
  =
  \hat k_\delta^+(x_N)\cdots\hat k_\delta^+(x_1)
  =
  \sum_{S\subset\{1,\ldots,N\}}
  \delta^{|S|}
  \prod_{j=N}^{1} \hat d_j(S),
  \label{eq:N-gain-subset-expansion}
\end{equation}
where
\begin{equation}
  \hat d_j(S)=
  \begin{cases}
    \hat c(x_j), & j\in S,\\
    \hat c^\dagger(x_j), & j\notin S .
  \end{cases}
  \label{eq:N-gain-subset-operator}
\end{equation}
The leading exact word is therefore
\begin{equation}
  \hat K_{+,N}(0)
  =
  \hat c^\dagger(x_N)\cdots\hat c^\dagger(x_1).
  \label{eq:N-gain-leading-word}
\end{equation}
If \(X\) has rank \(r=N\) and \(X=UR\), then
\begin{equation}
  \hat K_{+,N}(0)
  =
  \det(R)\,
  \hat c^\dagger(u_N)\cdots\hat c^\dagger(u_1),
  \qquad
  P_X=UU^\dagger .
  \label{eq:N-gain-qr-endpoint}
\end{equation}
The branch fills the endpoint subspace \(\operatorname{im}P_X\), the Born
factor is
\begin{equation}
  p_{+,X}(C)=\det\!\left[X^\dagger(\I-C)X\right],
  \label{eq:N-gain-born}
\end{equation}
and the active particle transfer is
\begin{equation}
  T_{\rm act}^{+X}=T_{\rm even}(\I-P_X).
  \label{eq:N-gain-active-transfer}
\end{equation}
The finite-regulator one-body product has \(N\) reciprocal zero/pole pairs
after the finite change of endpoint basis \(X\mapsto U\).  Rank deficiency of
\(X\), or singularity of the empty Gram \(X^\dagger(\I-C)X\), means that the
leading exact branch has zero wedge amplitude or zero Born probability.

\paragraph*{4. One loss.}
The loss branch is the particle-hole dual:
\begin{equation}
  \hat k_\delta^-(y)
  =
  \hat c(y)+\delta\hat c^\dagger(y)
  =
  \hat c(y)+O(\delta).
  \label{eq:one-loss-delta-limit}
\end{equation}
The exact map is the interior product \(\iota_y\).  Its Born factor is
\begin{equation}
  p_{-,y}(C)
  =
  \Tr\!\left[\hat c(y)\rho\hat c^\dagger(y)\right]
  =
  y^\dagger C y .
  \label{eq:one-loss-delta-born}
\end{equation}
If \(y=\|y\|\hat y\), the active finite chart is most naturally the hole
chart:
\begin{equation}
  T_{{\rm act},h}^{-y}
  =
  T_{{\rm even},h}(\I-P_{\hat y}^{\,T}).
  \label{eq:one-loss-active-transfer}
\end{equation}
The particle chart sees the reciprocal endpoint pole.  The endpoint
occupation label is empty rather than filled.

\paragraph*{5. Two losses.}
For two losses,
\begin{align}
  \hat k_\delta^-(y_2)\hat k_\delta^-(y_1)
  &=
  \hat c(y_2)\hat c(y_1)
  +\delta\!\left[
    \hat c(y_2)\hat c^\dagger(y_1)
    +\hat c^\dagger(y_2)\hat c(y_1)
  \right]
  +\delta^2\hat c^\dagger(y_2)\hat c^\dagger(y_1)
  \nonumber\\
  &=
  \hat c(y_2)\hat c(y_1)+O(\delta).
  \label{eq:two-loss-regulated-expansion}
\end{align}
The exact endpoint is a two-dimensional empty subspace if
\(y_1\wedge y_2\ne0\).  With \(Y=(y_1,y_2)\),
\begin{equation}
  p_{--,Y}(C)
  =
  \det
  \begin{pmatrix}
    y_1^\dagger C y_1 & y_1^\dagger C y_2\\
    y_2^\dagger C y_1 & y_2^\dagger C y_2
  \end{pmatrix}.
  \label{eq:two-loss-born-determinant}
\end{equation}
After \(Y=VR\), the endpoint projector is \(P_Y=VV^\dagger\), the active
hole transfer is \(T_{{\rm even},h}(\I-P_Y^T)\), and there are two
reciprocal regulator pairs.  Rank deficiency or a singular occupied Gram
means that the exact two-loss branch is redundant or blocked.

\paragraph*{6. \(M\) losses.}
For \(M\) losses,
\begin{equation}
  \hat K_{-,M}(\delta)
  =
  \hat k_\delta^-(y_M)\cdots\hat k_\delta^-(y_1)
  =
  \sum_{S\subset\{1,\ldots,M\}}
  \delta^{|S|}
  \prod_{j=M}^{1} \hat e_j(S),
  \label{eq:M-loss-subset-expansion}
\end{equation}
with
\begin{equation}
  \hat e_j(S)=
  \begin{cases}
    \hat c^\dagger(y_j), & j\in S,\\
    \hat c(y_j), & j\notin S .
  \end{cases}
  \label{eq:M-loss-subset-operator}
\end{equation}
Thus
\begin{equation}
  \hat K_{-,M}(0)=\hat c(y_M)\cdots\hat c(y_1).
  \label{eq:M-loss-leading-word}
\end{equation}
For independent \(Y=VR\), this is the oriented interior product on
\(\operatorname{im}P_Y\), multiplied by the determinant of the change of
basis.  The Born factor and active hole transfer are
\begin{equation}
  p_{-,Y}(C)=\det\!\left[Y^\dagger C Y\right],
  \qquad
  T_{{\rm act},h}^{-Y}=T_{{\rm even},h}(\I-P_Y^T).
  \label{eq:M-loss-born-active}
\end{equation}
The regulator produces \(M\) reciprocal endpoint pairs, now labelled as empty
sectors.

\paragraph*{7. One gain and one loss.}
There are two inequivalent orders.  If the gain occurs first and the loss
later, the operator product is
\begin{align}
  \hat k_\delta^-(\beta)\hat k_\delta^+(\alpha)
  &=
  \hat c(\beta)\hat c^\dagger(\alpha)
  +\delta\!\left[
    \hat c(\beta)\hat c(\alpha)
    +\hat c^\dagger(\beta)\hat c^\dagger(\alpha)
  \right]
  +\delta^2\hat c^\dagger(\beta)\hat c(\alpha)
  \nonumber\\
  &=
  \langle\beta|\alpha\rangle
  -\hat c^\dagger(\alpha)\hat c(\beta)
  +O(\delta).
  \label{eq:causal-pair-regulated-expansion}
\end{align}
With an even segment \(M\) between the gain and the loss, replace
\(\alpha\) by \(M\alpha\) in the contraction; the causal Gram is the scalar
\(\widetilde G=\langle\beta|M|\alpha\rangle\).  If
\(\widetilde G\ne0\), this is the one-pair version of
Eq.~\eqref{eq:balanced-block-master-formula}: a closed neutral block with
scalar \(\lambda=\widetilde G\) and a dressed rank-one subtraction in the
active graph chart.  For \(\beta=\alpha\) and \(M=\I\) it reduces to
\(\I-\hat n_\alpha\), the empty projector.  The Born factor for the exact
two-step branch is
\begin{equation}
  p_{-+}^{(0)}
  =
  \bigl[\alpha^\dagger(\I-C)\alpha\bigr]
  \bigl[\beta^\dagger C\beta\bigr]
  +
  \left|\beta^\dagger(\I-C)\alpha\right|^2 .
  \label{eq:causal-pair-born}
\end{equation}

If the loss occurs first and the gain later, then
\begin{align}
  \hat k_\delta^+(\alpha)\hat k_\delta^-(\beta)
  &=
  \hat c^\dagger(\alpha)\hat c(\beta)
  +\delta\!\left[
    \hat c^\dagger(\alpha)\hat c^\dagger(\beta)
    +\hat c(\alpha)\hat c(\beta)
  \right]
  +\delta^2\hat c(\alpha)\hat c^\dagger(\beta)
  \nonumber\\
  &=
  \hat c^\dagger(\alpha)\hat c(\beta)+O(\delta).
  \label{eq:anticausal-pair-regulated-expansion}
\end{align}
This is neutral in charge but not a scalar multiple of an invertible even
Gaussian transfer.  It is an exterior--interior map; for the same normalized
mode it is \(\hat n_\alpha\), the occupied projector.  Its Born factor is
\begin{equation}
  p_{+-}^{(0)}
  =
  \bigl[\beta^\dagger C\beta\bigr]
  \bigl[\alpha^\dagger(\I-C)\alpha\bigr]
  +
  \left|\alpha^\dagger C\beta\right|^2 .
  \label{eq:anticausal-pair-born}
\end{equation}
The two orders coincide only in special cases.  The regulator has not made
gain and loss commute; it has made their finite-\(\delta\) composition an
ordinary matrix product whose endpoint remembers the order.

\paragraph*{8. \(N\) gains and \(M\) losses.}
For a general chronological word, let
\(\sigma_j=+\) denote a gain and \(\sigma_j=-\) denote a loss.  Define
\begin{equation}
  d_j=
  \begin{cases}
    \hat c^\dagger(z_j), & \sigma_j=+,\\
    \hat c(z_j), & \sigma_j=-,
  \end{cases}
  \qquad
  \bar d_j=
  \begin{cases}
    \hat c(z_j), & \sigma_j=+,\\
    \hat c^\dagger(z_j), & \sigma_j=- .
  \end{cases}
  \label{eq:general-word-d-and-dbar}
\end{equation}
The regulated word is exactly
\begin{equation}
  \hat K_{\bm\sigma}(\delta)
  =
  \prod_{j=r}^{1}\left(d_j+\delta\bar d_j\right)
  =
  \sum_{S\subset\{1,\ldots,r\}}
  \delta^{|S|}
  \prod_{j=r}^{1} d_j(S),
  \qquad
  d_j(S)=
  \begin{cases}
    \bar d_j, & j\in S,\\
    d_j, & j\notin S .
  \end{cases}
  \label{eq:general-gain-loss-regulated-expansion}
\end{equation}
Therefore
\begin{equation}
  \hat K_{\bm\sigma}(\delta)
  =
  \hat K_{\bm\sigma}(0)+O(\delta),
  \qquad
  \hat K_{\bm\sigma}(0)=d_r\cdots d_1 .
  \label{eq:general-gain-loss-leading-word}
\end{equation}
For a \(U(1)\)-invariant incoming state, a single flipped branch changes the
net charge by two, so the Born probability satisfies
\begin{equation}
  \Tr[\hat K_{\bm\sigma}(\delta)\rho\hat K_{\bm\sigma}(\delta)^\dagger]
  =
  \Tr[\hat K_{\bm\sigma}(0)\rho\hat K_{\bm\sigma}(0)^\dagger]
  +O(\delta^2),
  \label{eq:general-regulated-born-convergence}
\end{equation}
unless the leading exact probability vanishes, in which case the limit is
nonuniform and the exact branch is blocked.

The leading word \(d_r\cdots d_1\) is then reduced by the Clifford relation
\(\iota_u\varepsilon_v+\varepsilon_v\iota_u=\langle u|v\rangle\).  Equivalently,
build the ordered Pfaffian matrix of
Eq.~\eqref{eq:mixed-word-born-pfaffian}.  For balanced causal blocks with
\(N=M\) and nonsingular causal Gram \(\widetilde G\),
\begin{equation}
  \hat K_{\bm\sigma}(0)
  =
  \lambda\,\hat K_T,
  \qquad
  \lambda=\pm\det\widetilde G,
  \qquad
  \frac{T}{\lambda}
  =
  M_{bb}-W_{fg}\widetilde G^{-1}W_l^\dagger .
  \label{eq:general-balanced-master-limit}
\end{equation}
More generally, let \(r_c\) be the rank of the causal contraction problem
after the chronological word has been cut into the blocks specified by
Eq.~\eqref{eq:running-charge-definition}.  The reduced word then has
\(N-r_c\) filled endpoint orbitals and \(M-r_c\) empty endpoint
orbitals/interior maps.  Thus even when \(N=M\), an anticausal or
rank-deficient ordering need not reduce to a matrix alone.  In all cases the
finite-\(\delta\) one-body matrices compose literally; the limit is not the
raw matrix product but the triple
\begin{equation}
  \left(
  \hbox{scalar Pfaffian/determinant},
  \quad
  \hbox{finite graph transfer},
  \quad
  \hbox{filled/empty endpoint sectors}
  \right).
  \label{eq:general-regulated-limit-data}
\end{equation}

\iffalse
The preceding product formula is not yet the exact gain/loss limit.  For
every \(\varepsilon>0\) it is literally ordinary Nambu, or equivalently
Majorana, matrix multiplication.  The exact operation is obtained only after
expanding the product in \(\varepsilon\) and separating three kinds of data:
finite active singular values, regulator zero/pole slopes, and the endpoint
occupation constraints carried by the exact word.  The following five cases
are the useful mechanical dictionary.

\paragraph*{1. One gain.}
Take a single normalized mode \(\chi\) and the regulated gain
\begin{equation}
  \hat K_\varepsilon^+
  =
  \hat\chi^\dagger+\varepsilon\hat\chi
  \xrightarrow{\varepsilon\to0}
  \hat\chi^\dagger .
  \label{eq:one-gain-regulated-limit}
\end{equation}
In the Nambu basis its one-body representative is
Eq.~\eqref{eq:nambu-regulated-gain-matrix}.  If the neighboring even pieces
are number conserving,
\(\mathbb A_L=\operatorname{diag}(A_{L,h},A_{L,p})\) and
\(\mathbb A_R=\operatorname{diag}(A_{R,h},A_{R,p})\), then direct
multiplication gives
\begin{equation}
  \mathbb A_L\mathbb R_+(\chi,\varepsilon)\mathbb A_R
  =
  \begin{pmatrix}
    -A_{L,h}(\I-P)A_{R,h}
    &
    \varepsilon^{-1}A_{L,h}\chi\chi^T A_{R,p}
    \\
    \varepsilon A_{L,p}\chi^*\chi^\dagger A_{R,h}
    &
    -A_{L,p}(\I-P^T)A_{R,p}
  \end{pmatrix}.
  \label{eq:one-gain-sandwiched-limit}
\end{equation}
Thus the ordinary entrywise limit is not a full transfer matrix: the
particle-hole block has a pole, the reciprocal hole-particle block has a
zero, and the diagonal fixed-chart blocks become the reduced maps
\begin{equation}
  A_{L,h}A_{R,h}\mapsto A_{L,h}(\I-P)A_{R,h},
  \qquad
  A_{L,p}A_{R,p}\mapsto A_{L,p}(\I-P^T)A_{R,p}.
  \label{eq:one-gain-reduced-diagonal-blocks}
\end{equation}
The exact branch is instead the finite exterior operation
\(\hat\chi^\dagger\).  It fixes one filled endpoint mode and has Born weight
\begin{equation}
  p_+(C)=\Tr(\hat\chi^\dagger\rho\hat\chi)=1-\chi^\dagger C\chi .
  \label{eq:one-gain-born-limit}
\end{equation}
Consequently a one-gain regulator contributes one zero/pole singular pair
with slopes \(\kappa=\pm1\), plus a finite reduced transfer on the complement
of \(\chi\).  The divergent pair is physical endpoint bookkeeping, not a
finite Lyapunov gap.

\paragraph*{2. \(r\) gains.}
Now consider a same-type gain block with \(r\) exact creations separated by
number-conserving even Gaussian pieces.  Move all creations to a common edge
by defining the transported modes through
\begin{equation}
  \hat E_r\cdots\hat E_j\,\hat\chi_j^\dagger\,
  (\hat E_r\cdots\hat E_j)^{-1}
  =
  \hat c^\dagger(\tilde\chi_j),
  \qquad
  X=(\tilde\chi_1,\ldots,\tilde\chi_r).
  \label{eq:many-gain-transported-modes}
\end{equation}
The exact branch is the exterior product
\(\hat c^\dagger(\tilde\chi_r)\cdots\hat c^\dagger(\tilde\chi_1)\) times the
remaining even Gaussian map.  Its sign depends on the order, but for a
same-type block all order dependence is the fermionic orientation of the
wedge product.  If \(X\) has rank \(r_+\), the branch fixes the
\(r_+\)-dimensional filled subspace
\(\mathcal X=\operatorname{span}X\).  For independent modes,
\begin{equation}
  P_X=X(X^\dagger X)^{-1}X^\dagger
  \label{eq:many-gain-span-projector}
\end{equation}
is the basis-independent Hermitian projector onto that span; for rank
deficient \(X\) the inverse is replaced by the Moore--Penrose inverse and the
antisymmetrized exact branch has rank \(r_+<r\), or vanishes if a repeated
constraint is imposed exactly.

The Born weight of the same gain block is the determinant of the empty-space
Gram matrix,
\begin{equation}
  p_{+,\mathcal X}(C)
  =
  \det\!\left[X^\dagger(\I-C)X\right],
  \label{eq:many-gain-born-limit}
\end{equation}
with the determinant taken on the independent constrained subspace.  The
corresponding covariance update is Eq.~\eqref{eq:many-gain-covariance}.
In the fixed particle/hole chart the finite active block is the transfer
compressed to the complement of \(\mathcal X\),
\begin{equation}
  A_h\mapsto A_h(\I-P_X),
  \qquad
  A_p\mapsto A_p(\I-P_X^T),
  \label{eq:many-gain-reduced-transfer}
\end{equation}
again up to left/right placement conventions set by where the block is
collected.  In the regulated Nambu product this same information appears as
\(r_+\) reciprocal zero/pole pairs and a finite reduced singular spectrum.
Non-orthogonal gain wavefunctions cause no new algebraic rule; they only make
the Gram matrix in Eq.~\eqref{eq:many-gain-born-limit} non-diagonal and can
make the branch ill-conditioned or blocked when the Gram determinant is small
or zero.

\paragraph*{3. One loss.}
The single-loss limit is the particle-hole dual of the previous case.  With
\begin{equation}
  \hat K_\varepsilon^-
  =
  \hat\chi+\varepsilon\hat\chi^\dagger
  \xrightarrow{\varepsilon\to0}
  \hat\chi ,
  \label{eq:one-loss-regulated-limit}
\end{equation}
Eq.~\eqref{eq:nambu-regulated-loss-matrix} gives
\begin{equation}
  \mathbb A_L\mathbb R_-(\chi,\varepsilon)\mathbb A_R
  =
  \begin{pmatrix}
    -A_{L,h}(\I-P)A_{R,h}
    &
    \varepsilon A_{L,h}\chi\chi^T A_{R,p}
    \\
    \varepsilon^{-1}A_{L,p}\chi^*\chi^\dagger A_{R,h}
    &
    -A_{L,p}(\I-P^T)A_{R,p}
  \end{pmatrix}.
  \label{eq:one-loss-sandwiched-limit}
\end{equation}
The finite diagonal compression is the same as in
Eq.~\eqref{eq:one-gain-reduced-diagonal-blocks}; the divergent off-diagonal
block is the opposite one.  The endpoint label is now empty rather than
filled, and
\begin{equation}
  p_-(C)=\Tr(\hat\chi\rho\hat\chi^\dagger)=\chi^\dagger C\chi .
  \label{eq:one-loss-born-limit}
\end{equation}
Thus one loss also contributes one reciprocal regulator pair with
\(\kappa=\pm1\), but it represents an empty endpoint constraint.

\paragraph*{4. \(r\) losses.}
For a same-type loss block, move all annihilators to a common edge by
\begin{equation}
  (\hat E_{j-1}\cdots\hat E_0)^{-1}\,\hat\chi_j\,
  (\hat E_{j-1}\cdots\hat E_0)
  =
  \hat c(\tilde\eta_j),
  \qquad
  Y=(\tilde\eta_1,\ldots,\tilde\eta_r).
  \label{eq:many-loss-transported-modes}
\end{equation}
The exact branch is an ordered interior product by the modes in \(Y\).  If
\(Y\) has independent rank \(r_-\), it fixes an \(r_-\)-dimensional empty
endpoint subspace \(\mathcal Y=\operatorname{span}Y\).  Its Born weight is
the occupied-space Gram determinant
\begin{equation}
  p_{-,\mathcal Y}(C)
  =
  \det\!\left[Y^\dagger C Y\right],
  \label{eq:many-loss-born-limit}
\end{equation}
and the covariance update is Eq.~\eqref{eq:many-loss-covariance}.  The
finite fixed-chart transfer is compressed to the complement of the empty
constraint,
\begin{equation}
  A_h\mapsto A_h(\I-P_Y),
  \qquad
  A_p\mapsto A_p(\I-P_Y^T),
  \qquad
  P_Y=Y(Y^\dagger Y)^+Y^\dagger .
  \label{eq:many-loss-reduced-transfer}
\end{equation}
As for gains, non-orthogonal wavefunctions are handled by the full Gram
matrix.  The regulated Nambu product produces \(r_-\) zero/pole pairs whose
endpoint label is empty and whose finite complement is the reduced transfer
in Eq.~\eqref{eq:many-loss-reduced-transfer}.

\paragraph*{5. General non-commuting gain/loss words.}
For mixed words there is no simplification to separate gain and loss spans
before the ordered product is evaluated.  Gain and loss do not commute, and
even for two non-orthogonal modes one has
\begin{equation}
  \hat\beta\,\hat\alpha^\dagger
  =
  \langle\beta|\alpha\rangle
  -
  \hat\alpha^\dagger\hat\beta ,
  \label{eq:mixed-gain-loss-contraction-order}
\end{equation}
so changing the order changes both the scalar contraction and the residual
Clifford word.  The robust rule is therefore to keep the finite-regulator
Nambu product ordered and only then take the Laurent limit.  Write every
regulated odd reflection as
\begin{equation}
  \mathbb R_j(\varepsilon)
  =
  \varepsilon^{-1}\mathbb R_j^{[-1]}
  +
  \mathbb R_j^{[0]}
  +
  \varepsilon\mathbb R_j^{[+1]} .
  \label{eq:regulated-reflection-laurent}
\end{equation}
For gain,
\begin{align}
  \mathbb R_+^{[-1]}
  &=
  \begin{pmatrix}0&\chi\chi^T\\0&0\end{pmatrix},
  &
  \mathbb R_+^{[0]}
  &=
  -
  \begin{pmatrix}\I-P&0\\0&\I-P^T\end{pmatrix},
  &
  \mathbb R_+^{[+1]}
  &=
  \begin{pmatrix}0&0\\\chi^*\chi^\dagger&0\end{pmatrix},
  \label{eq:gain-reflection-laurent-coeffs}
\end{align}
while for loss the off-diagonal Laurent coefficients are interchanged,
\begin{align}
  \mathbb R_-^{[-1]}
  &=
  \begin{pmatrix}0&0\\\chi^*\chi^\dagger&0\end{pmatrix},
  &
  \mathbb R_-^{[0]}
  &=
  -
  \begin{pmatrix}\I-P&0\\0&\I-P^T\end{pmatrix},
  &
  \mathbb R_-^{[+1]}
  &=
  \begin{pmatrix}0&\chi\chi^T\\0&0\end{pmatrix}.
  \label{eq:loss-reflection-laurent-coeffs}
\end{align}
Substitution into Eq.~\eqref{eq:arbitrary-regulated-nambu-product} gives the
finite Laurent series
\begin{equation}
  \mathbb T_{\bxi,\varepsilon}
  =
  \sum_{q=-m}^{m}\varepsilon^q
  \mathbb T_\bxi^{[q]},
  \label{eq:regulated-word-laurent-series}
\end{equation}
with coefficients
\begin{equation}
  \mathbb T_\bxi^{[q]}
  =
  \sum_{\substack{\ell_1,\ldots,\ell_m\in\{-1,0,+1\}\\
  \ell_1+\cdots+\ell_m=q}}
  \mathbb A_0\mathbb R_1^{[\ell_1]}\mathbb A_1\cdots
  \mathbb R_m^{[\ell_m]}\mathbb A_m .
  \label{eq:regulated-word-laurent-coefficients}
\end{equation}
Equation~\eqref{eq:regulated-word-laurent-coefficients} is the mechanical
answer for arbitrary non-commuting gain/loss strings.  The leading pole is
not determined by the number of gain/loss events alone: nilpotent products,
nonzero contractions, vanishing overlaps, and rank-deficient constrained
subspaces can all change the leading Laurent coefficient.

Let
\begin{equation}
  q_{\min}=\min\{q:\mathbb T_\bxi^{[q]}\ne0\},
  \qquad
  q_{\max}=\max\{q:\mathbb T_\bxi^{[q]}\ne0\}.
  \label{eq:word-laurent-extreme-powers}
\end{equation}
If \(q_{\min}<0\), the entrywise matrix diverges, but its projective leading
flag is well defined:
\begin{equation}
  \varepsilon^{-q_{\min}}\mathbb T_{\bxi,\varepsilon}
  \xrightarrow{\varepsilon\to0}
  \mathbb T_\bxi^{[q_{\min}]} .
  \label{eq:leading-laurent-flag}
\end{equation}
This coefficient is not an invertible Gaussian transfer matrix.  Since the
finite-regulator product obeys
\(\mathbb T_{\bxi,\varepsilon}\Omega\mathbb T_{\bxi,\varepsilon}^T=\Omega\),
the Laurent coefficients satisfy
\begin{equation}
  \sum_{q+q'=r}
  \mathbb T_\bxi^{[q]}\Omega
  \bigl(\mathbb T_\bxi^{[q']}\bigr)^T
  =
  \delta_{r,0}\Omega .
  \label{eq:laurent-orthogonality-identities}
\end{equation}
The extreme zero/pole coefficients are therefore isotropic flags.  They are
the Nambu/Majorana representation of exact occupation constraints, not
finite transfer blocks.

The finite spectrum is extracted from singular-value asymptotics,
\begin{equation}
  \log s_a(\varepsilon)
  =
  \kappa_a\log\frac{1}{\varepsilon}
  +
  \rho_a
  +
  o(1),
  \qquad
  a=1,\ldots,2N .
  \label{eq:regulated-singular-value-asymptotics}
\end{equation}
CAR preservation pairs the regulator slopes \(\kappa_a\) with opposite
slopes.  Nonzero \(\kappa_a\) are the endpoint or gauge sectors.  The
regulator-stable offsets \(\rho_a\) with \(\kappa_a=0\), together with the
filled/empty endpoint labels obtained from the exact ordered word, are the
complete singular data used in the Choi transfer convention.

Equivalently, if \(V_0(\varepsilon)\) and \(U_0(\varepsilon)\) are the right
and left singular subspaces with \(\kappa_a=0\), the finite active transfer is
represented, up to finite chart rotations, by
\begin{equation}
  U_0(\varepsilon)^\dagger
  \mathbb T_{\bxi,\varepsilon}
  V_0(\varepsilon)
  \xrightarrow{\varepsilon\to0}
  \hbox{finite singular values }e^{\rho_a}.
  \label{eq:finite-sector-compression}
\end{equation}
There is generally no unique full finite matrix limit.  There is a unique
finite singular spectrum only after the regulator slopes and the scalar
normalization have been separated.

A one-mode mixed word shows how a regulator returns an even projective
measurement.  With
\(\hat K^+_\varepsilon=\hat\chi^\dagger+\varepsilon\hat\chi\) and
\(\hat K^-_\varepsilon=\hat\chi+\varepsilon\hat\chi^\dagger\),
\begin{align}
  \hat K^-_\varepsilon\hat K^+_\varepsilon
  &=
  \hat\chi\hat\chi^\dagger
  +
  \varepsilon^2\hat\chi^\dagger\hat\chi
  =
  \Pi_0+\varepsilon^2\Pi_1,
  \label{eq:loss-gain-product-empty-projector}\\
  \hat K^+_\varepsilon\hat K^-_\varepsilon
  &=
  \hat\chi^\dagger\hat\chi
  +
  \varepsilon^2\hat\chi\hat\chi^\dagger
  =
  \Pi_1+\varepsilon^2\Pi_0 .
  \label{eq:gain-loss-product-occupied-projector}
\end{align}
Thus loss after gain limits to the empty projector, while gain after loss
limits to the occupied projector.  In the same Nambu pair,
\begin{equation}
  \mathbb R_+(\varepsilon)\mathbb R_-(\varepsilon)
  =
  \begin{pmatrix}\varepsilon^{-2}&0\\0&\varepsilon^2\end{pmatrix},
  \qquad
  \mathbb R_-(\varepsilon)\mathbb R_+(\varepsilon)
  =
  \begin{pmatrix}\varepsilon^2&0\\0&\varepsilon^{-2}\end{pmatrix}.
  \label{eq:odd-pair-projector-transfer-limit}
\end{equation}
The product of two odd regulated transfers has become the zero/pole transfer
of an even occupation measurement.  For non-orthogonal modes this statement
is not diagonal in the original mode labels; one must use the ordered
contraction/Pfaffian matrix of Eq.~\eqref{eq:mixed-word-born-pfaffian}, or
equivalently the ordered Laurent coefficients above.  The exact limit is
taken after the ordered regulated product has been formed.

The Majorana version is the same construction after a fixed change of basis.
Let \(\gamma=S\Psi\), so the CAR form becomes the identity and
\(R_{\bxi,\varepsilon}=S\mathbb T_{\bxi,\varepsilon}S^{-1}\).  A linear form
\(\hat\ell(w)=u^T\gamma\) with \(u^Tu\ne0\) gives
\begin{equation}
  R[u]=-
  \I+\frac{2uu^T}{u^Tu},
  \qquad
  R[u]R[u]^T=
  \I,
  \label{eq:majorana-reflection-arbitrary-u}
\end{equation}
and the arbitrary word is the corresponding ordered product of even
\(SO(2N,\mathbb C)\) matrices and odd reflections.  The exact gain/loss limit
is the limit in which selected Nambu or Majorana vectors approach the null
cone.  The transfer matrix may then diverge or cease to be a graph, but the
Choi covariance has a finite limit.
\fi

The Born weight of the same regulated product is not determined by the
projective one-body matrix \(\mathbb T_{\bxi,\varepsilon}\), or by its
Majorana-basis representative, alone.  For one branch and an even Gaussian
density matrix, define
\begin{equation}
  M_{ab}=\Tr(\rho\,\gamma_a\gamma_b).
  \label{eq:majorana-two-point-M}
\end{equation}
Equation~\eqref{eq:regulated-branch-pin-factorization} gives
\begin{equation}
  p_j(\varepsilon)
  =
  \Tr[\hat K_j(\varepsilon)\rho\hat K_j(\varepsilon)^\dagger]
  =
  |c_j(\varepsilon)|^2\,
  u_j(\varepsilon)^\dagger M u_j(\varepsilon).
  \label{eq:single-regulated-majorana-born}
\end{equation}
For a same-mode number-conserving state and
\(\hat K=a\hat\chi+b\hat\chi^\dagger\), this reduces to
\begin{equation}
  p_{a,b}(C)=|a|^2 n_\chi+|b|^2(1-n_\chi),
  \qquad
  n_\chi=\bra{\chi}C\ket{\chi} .
  \label{eq:single-regulated-u1-born}
\end{equation}
Therefore
\begin{equation}
  p^+_\varepsilon(C)=1-n_\chi+\varepsilon^2 n_\chi,
  \qquad
  p^-_\varepsilon(C)=n_\chi+\varepsilon^2(1-n_\chi),
  \label{eq:regulated-gain-loss-born-correct}
\end{equation}
which tends to the exact gain/loss probabilities in
Eq.~\eqref{eq:exact-gain-loss-born}.

For an ordered string of regulated odd insertions,
\begin{equation}
  \hat K_{\bxi,\varepsilon}
  =
  \left[\prod_{j=1}^r c_j(\varepsilon)\right]
  \hat u_r(\varepsilon)\cdots\hat u_1(\varepsilon) .
  \label{eq:regulated-string-factorization}
\end{equation}
Let
\begin{equation}
  L_1,\ldots,L_{2r}
  =
  \hat u_1^\dagger,\ldots,\hat u_r^\dagger,
  \hat u_r,\ldots,\hat u_1
  \label{eq:regulated-product-L-list}
\end{equation}
be the ordered list entering
\(\hat K_{\bxi,\varepsilon}^\dagger\hat K_{\bxi,\varepsilon}\).  With
\begin{equation}
  W_{\mu\nu}^{(\varepsilon)}=
  \begin{cases}
    \langle L_\mu L_\nu\rangle, & \mu<\nu,\\
    -\langle L_\nu L_\mu\rangle, & \mu>\nu,\\
    0, & \mu=\nu,
  \end{cases}
  \label{eq:regulated-product-wick-matrix}
\end{equation}
Wick's theorem gives
\begin{equation}
  p_{\bxi,\varepsilon}
  =
  \Tr(\hat K_{\bxi,\varepsilon}\rho
  \hat K_{\bxi,\varepsilon}^\dagger)
  =
  \left|\prod_{j=1}^r c_j(\varepsilon)\right|^2
  \operatorname{Pf}W^{(\varepsilon)} .
  \label{eq:regulated-product-born-pfaffian}
\end{equation}
Even Gaussian segments are included by inserting their own Bogoliubov transfer
matrices into Eq.~\eqref{eq:arbitrary-regulated-nambu-product} and by keeping
their determinant/scalar normalizations in the same scalar cocycle.  Thus the
projective matrix product and the scalar/Pfaffian normalization are two pieces
of one trajectory amplitude, not interchangeable quantities.

For fixed finite string length, the exact word is recovered as
\begin{equation}
  \hat K_{\bxi,\varepsilon}=\hat K_{\bxi,0}+O(\varepsilon),
  \label{eq:regulated-word-limit}
\end{equation}
where \(\hat K_{\bxi,0}\) is the chronological word of exact
\(\hat\chi^\dagger\)'s and \(\hat\chi\)'s.  If the incoming Gaussian state is
\(U(1)\)-invariant, terms with a single flipped branch carry the wrong charge
and do not interfere with the exact word, so generically
\begin{equation}
  p_{\bxi,\varepsilon}=p_{\bxi,0}+O(\varepsilon^2).
  \label{eq:regulated-born-limit}
\end{equation}
Without \(U(1)\) symmetry the convergence is still to the exact word, but the
leading correction can be \(O(\varepsilon)\).  If \(p_{\bxi,0}=0\), the limit
is nonuniform: after normalization the regulated state is controlled by the
first nonzero \(\varepsilon\)-suppressed word, not by the blocked exact branch.
Such events should be treated as censored exact gain/loss branches.

Finally, let \(s_a(\varepsilon)\) be singular values of the projective product
\(\mathbb T_{\bxi,\varepsilon}\).  Regulator directions are detected by
\begin{equation}
  \kappa_a=
  \lim_{\varepsilon\to0}
  \frac{\log s_a(\varepsilon)}{\log(1/\varepsilon)} .
  \label{eq:regulator-slope-kappa}
\end{equation}
Nonzero \(\kappa_a\) are endpoint/gauge sectors.  The finite transfer gaps
are the regulator-stable singular data, or Fock-space combinations whose
\(\log\varepsilon\) dependence cancels after the scalar normalization in
Eq.~\eqref{eq:regulated-product-born-pfaffian} is included.  This is why a
regulated Majorana basis simplifies composition but does not remove endpoint
bookkeeping or Born-weight bookkeeping.

\subsection{Dressed Projectors from the Single-Copy Nambu Regulator}
\label{sec:nambu-regulator-dressed-projector}

The dressed projector in the Choi-contraction calculation can be obtained
without passing to the doubled superoperator representation.  The only
essential point is the order of limits: first compose the finite-regulator
Nambu word, or equivalently transport the regulated odd linear form through
the later even factor; only then take the singular endpoint limit.

In this subsection it is cleaner to remove the antihomomorphism ambiguity by
using the circuit-order matrix
\begin{equation}
  L[\hat V]\equiv \mathbb T[\hat V]^T,
  \qquad
  L[\hat V_2\hat V_1]=L[\hat V_2]L[\hat V_1].
  \label{eq:circuit-order-L-convention}
\end{equation}
This is only a local convention; elsewhere the note keeps the
\(\mathbb T\)-convention of Eq.~\eqref{eq:nambu-antihomomorphism}.  For a
gain in mode \(\chi\), define the null Nambu covectors
\begin{equation}
  w_\chi=
  \begin{pmatrix}0\\ \chi\end{pmatrix},
  \qquad
  \bar u_\chi=
  \begin{pmatrix}\chi^*\\0\end{pmatrix},
  \qquad
  w_\chi^T\Omega \bar u_\chi=1,
  \label{eq:single-copy-null-gain-pair}
\end{equation}
and the normalized non-null regulator
\begin{equation}
  \zeta_\varepsilon
  =
  \varepsilon^{-1/2}w_\chi
  +
  \varepsilon^{1/2}\bar u_\chi,
  \qquad
  \zeta_\varepsilon^T\Omega\zeta_\varepsilon=2 .
  \label{eq:single-copy-gain-regulator-zeta}
\end{equation}
The corresponding Pin element is
\(\hat g(\varepsilon)=\hat\ell(\zeta_\varepsilon)\), and its reflection is
\begin{equation}
  L_g(\varepsilon)
  =
  \zeta_\varepsilon\zeta_\varepsilon^T\Omega-\I,
  \qquad
  L_g(\varepsilon)\Omega L_g(\varepsilon)^T=\Omega .
  \label{eq:L-reflection-form}
\end{equation}
The projectively normalized operator used below,
\(\hat k_\varepsilon^+(\chi)=\hat c^\dagger(\chi)+\varepsilon\hat c(\chi)\),
is \(\sqrt{\varepsilon}\,\hat\ell(\zeta_\varepsilon)\).  The scalar
\(\sqrt{\varepsilon}\) belongs to the amplitude cocycle and is not visible in
the projective matrix \(L_g\).

Now sandwich the regulated gain by two invertible even Gaussian pieces,
\(\hat K(\varepsilon)=\hat V_2\hat g(\varepsilon)\hat V_1\), and let
\(L_i=L[\hat V_i]\).  Orthogonality gives
\begin{equation}
  L_K(\varepsilon)
  =
  L_2L_g(\varepsilon)L_1
  =
  a(\varepsilon)b(\varepsilon)^T\Omega-L_0,
  \qquad
  a=L_2\zeta_\varepsilon,\quad
  b=L_1^{-1}\zeta_\varepsilon,\quad
  L_0=L_2L_1 .
  \label{eq:single-insertion-rank-one-correction}
\end{equation}
This is an exact finite-\(\varepsilon\) identity in \(O(2N,\mathbb C)\).  Its
projectively rescaled matrix limit is
\begin{equation}
  \varepsilon L_K(\varepsilon)
  \xrightarrow{\varepsilon\to0}
  (L_2w_\chi)(L_1^{-1}w_\chi)^T\Omega ,
  \label{eq:single-insertion-rank-one-projective-limit}
\end{equation}
which has rank one.  That rank-one flag is not a valid transfer matrix on the
full graph.  The finite dressed projector is obtained only after moving to
the appropriate graph/covariance chart, or equivalently after transporting
the regulated null vector before taking the endpoint limit.

Let \(\hat V\) and \(\hat W\) be number-conserving even Gaussian operators
with particle transfer matrices \(T\) and \(S\), using the convention
\[
  T_p(\hat W\hat V)=T_p(\hat V)T_p(\hat W)=TS .
\]
For a one-particle wavefunction \(x\), write
\begin{equation}
  \hat c^\dagger(x)=\sum_i x_i\hat c_i^\dagger,
  \qquad
  \hat c(x)=\sum_i x_i^*\hat c_i .
  \label{eq:mode-linear-forms-for-dressed-projector}
\end{equation}
The even operator \(\hat W\) acts on creation and annihilation covectors as
\begin{equation}
  \hat W\,\hat c^\dagger(\chi)\,\hat W^{-1}
  =
  \hat c^\dagger(S^\dagger\chi),
  \qquad
  \hat W\,\hat c(\chi)\,\hat W^{-1}
  =
  \hat c(S^{-1}\chi).
  \label{eq:even-dresses-nambu-linear-forms}
\end{equation}
Equivalently, with \(\Psi=(c,c^\dagger)^T\) and
\(\hat\ell(w)=w^T\Psi\), the gain regulator has Nambu covector
\begin{equation}
  w_\varepsilon^+(\chi)
  =
  \begin{pmatrix}
    \varepsilon\chi^*\\
    \chi
  \end{pmatrix},
  \qquad
  \hat\ell[w_\varepsilon^+(\chi)]
  =
  \hat c^\dagger(\chi)+\varepsilon\hat c(\chi).
  \label{eq:nambu-covector-gain-regulator}
\end{equation}
Finite Nambu matrix composition transports this covector to
\begin{equation}
  \widetilde w_\varepsilon^+(\chi;S)
  =
  \begin{pmatrix}
    \varepsilon(S^{-1}\chi)^*\\
    S^\dagger\chi
  \end{pmatrix}.
  \label{eq:nambu-transported-gain-covector}
\end{equation}
Thus the finite Nambu regulator for a gain insertion,
\begin{equation}
  \hat k_\varepsilon^+(\chi)
  =
  \hat c^\dagger(\chi)+\varepsilon\hat c(\chi),
  \qquad
  \bigl[\hat k_\varepsilon^+(\chi)\bigr]^2=\varepsilon,
  \label{eq:single-copy-gain-regulator}
\end{equation}
gives the exact finite-regulator identity
\begin{align}
  \hat W\,\hat k_\varepsilon^+(\chi)\,\hat V
  &=
  \left[
    \hat c^\dagger(x)+\varepsilon\hat c(y)
  \right]\hat W\hat V,
  \label{eq:transported-gain-regulator}\\
  x&=S^\dagger\chi,
  \qquad
  y=S^{-1}\chi,
  \qquad
  y^\dagger x=1 .
  \label{eq:biorthogonal-dressed-gain-pair}
\end{align}
Equation~\eqref{eq:transported-gain-regulator} is the single-copy Nambu
version of the dressing rule.  After a nonunitary \(S\), the regulator partner
is not the Hermitian conjugate wavefunction of the dominant creation mode.
The pair \((x,y)\) is instead biorthogonal.  This is why freezing the
undressed projector \(P_\chi\) before composing with \(S\) gives the wrong
finite chart.

Since the projective limit is controlled by the creation component \(x\), put
\begin{equation}
  \hat x=\frac{x}{\|x\|},
  \qquad
  P_{\hat x}=\hat x\hat x^\dagger
  =
  \frac{S^\dagger\chi\chi^\dagger S}{\chi^\dagger SS^\dagger\chi}.
  \label{eq:dressed-gain-projector-nambu}
\end{equation}
The scalar \(\|x\|\) belongs to the branch normalization.  It affects the
Born weight but not the finite projective transfer block.  The exact gain
branch is
\begin{equation}
  \hat W\hat\chi^\dagger\hat V
  =
  \|x\|\,
  \hat c^\dagger(\hat x)\,\hat W\hat V .
  \label{eq:gain-after-transport-projective-limit}
\end{equation}
Using
\(\hat c^\dagger(\hat x)=\hat c^\dagger(\hat x)(\I-\hat n_{\hat x})\), fix
the particle-chart gauge of the even part by
\begin{equation}
  \I-\hat n_{\hat x}
  =
  \lim_{\mu\to\infty}e^{-\mu\hat n_{\hat x}},
  \qquad
  D_\mu(\hat x)=\I-P_{\hat x}+e^{-\mu}P_{\hat x}.
  \label{eq:dressed-gain-empty-gauge}
\end{equation}
At finite \(\mu\), this is an ordinary number-conserving Nambu transfer; its
particle block is multiplied by \(D_\mu(\hat x)\).  Therefore
\begin{equation}
  T_{\rm eff}^{+}
  =
  \lim_{\mu\to\infty}
  T_p(\hat W\hat V)\,D_\mu(\hat x)
  =
  TS\left(\I-P_{\hat x}\right)
  =
  TS\left[
    \I-
    \frac{S^\dagger\chi\chi^\dagger S}
         {\chi^\dagger SS^\dagger\chi}
  \right].
  \label{eq:nambu-derived-dressed-gain-transfer}
\end{equation}
This is exactly the dressed-projector formula obtained from the Choi
contraction calculation.  It is a single-copy Nambu result.  The apparent
extra projector dressing comes from transporting the \emph{regulated null
vector} through the later nonunitary even segment before collapsing the
endpoint sector.

The loss formula follows by the same argument in the hole chart.  Start from
\begin{equation}
  \hat k_\varepsilon^-(\chi)
  =
  \hat c(\chi)+\varepsilon\hat c^\dagger(\chi).
  \label{eq:single-copy-loss-regulator}
\end{equation}
In Nambu-covector form,
\begin{equation}
  w_\varepsilon^-(\chi)
  =
  \begin{pmatrix}
    \chi^*\\
    \varepsilon\chi
  \end{pmatrix}
  \quad\longmapsto\quad
  \widetilde w_\varepsilon^-(\chi;S)
  =
  \begin{pmatrix}
    (S^{-1}\chi)^*\\
    \varepsilon S^\dagger\chi
  \end{pmatrix}.
  \label{eq:nambu-transported-loss-covector}
\end{equation}
Equivalently, conjugating through \(\hat W\) gives
\begin{equation}
  \hat W\,\hat k_\varepsilon^-(\chi)\,\hat V
  =
  \left[
    \hat c(y)+\varepsilon\hat c^\dagger(x)
  \right]\hat W\hat V,
  \qquad
  y=S^{-1}\chi,\quad x=S^\dagger\chi .
  \label{eq:transported-loss-regulator}
\end{equation}
The dominant exact operator is now the annihilator in the normalized mode
\(\hat y=y/\|y\|\), with
\begin{equation}
  P_{\hat y}
  =
  \frac{S^{-1}\chi\chi^\dagger S^{-\dagger}}
       {\chi^\dagger S^{-\dagger}S^{-1}\chi}.
  \label{eq:dressed-loss-projector-nambu}
\end{equation}
Because loss is naturally reconstructed from the hole dictionary, the finite
hole transfer is
\begin{equation}
  T_{{\rm eff},h}^{-}
  =
  (TS)^{-T}\left(\I-P_{\hat y}^{\,T}\right).
  \label{eq:nambu-derived-dressed-loss-transfer}
\end{equation}
The homogeneous particle-side loss formula alone misses the endpoint source
created by the exact annihilation; the hole formula
Eq.~\eqref{eq:nambu-derived-dressed-loss-transfer} is the finite chart in
which the singular sector has been removed.

Equations~\eqref{eq:nambu-derived-dressed-gain-transfer} and
\eqref{eq:nambu-derived-dressed-loss-transfer} also explain the common
mistake.  If one takes the endpoint limit at the insertion before composing
with \(\hat W\), one writes a projector \(P_\chi\) in the wrong frame.  The
correct Nambu limit is
\begin{equation}
\begin{gathered}
  \text{multiply/transport at finite regulator}
  \longrightarrow
  \text{identify the dominant null vector}
  \\
  \longrightarrow
  \text{take the rank-deficient endpoint}.
\end{gathered}
\label{eq:correct-nambu-limit-order}
\end{equation}
For unitary \(S\), \(S^{-1}\chi=S^\dagger\chi\) and the distinction collapses;
for nonunitary monitored evolution it is essential.

\subsection{Neutral Choi Reference Versus Odd-Operator Regularization}
\label{sec:neutral-choi-doubled-regulator}

There is an important distinction between two regularizations that are easy to
conflate.  The Pin/Nambu regulator above regularizes the odd trajectory
amplitude \(\hat K_\bxi\) itself.  The contraction rule in the Choi notes,
however, regularizes the \emph{channel} or superoperator
\begin{equation}
  \mathcal J_K:\rho\mapsto \hat K\rho\hat K^\dagger
  \label{eq:jump-superoperator-definition}
\end{equation}
after vectorization.  These are different one-body charts.  They have the
same exact endpoint after the correct scalar/postselection limit is taken, but
their finite-regulator matrices do not have to look the same.

The clean Choi construction starts from a neutral maximally entangled
reference.  In a Majorana-neutral convention one may take, mode by mode,
\begin{equation}
  \ket{\Omega_M}_{LR}
  =
  2^{-N/2}\prod_{i=1}^N
  \left(1+c_{L,i}^\dagger c_{R,i}^\dagger\right)\ket{0},
  \label{eq:majorana-neutral-reference}
\end{equation}
for which a right-leg insertion can be moved to the left leg by a fixed
particle-hole mirror, up to the usual graded signs:
\begin{equation}
  c_{R,i}\ket{\Omega_M}_{LR}
  =
  -c_{L,i}^\dagger\ket{\Omega_M}_{LR},
  \qquad
  c_{R,i}^\dagger\ket{\Omega_M}_{LR}
  =
  c_{L,i}\ket{\Omega_M}_{LR}.
  \label{eq:neutral-reference-intertwiner}
\end{equation}
Thus, on the bare reference, left and right insertions are equivalent
descriptions.  After an operator has already acted, the same move becomes a
graded transpose/mirror of the operator order; it does not permit one to
commute a gain through later nonunitary factors.

For comparison with the Choi-contraction notes, one may also vectorize the channel
\(\mathcal J_K\).  Introduce ket and bra copies, labelled \(+\) and \(-\), so
that a gain jump is represented by the even bilinear
\begin{equation}
  J_{\rm gain}[\chi]
  =
  \hat\chi_+^\dagger\hat\chi_- ,
  \qquad
  J_{\rm loss}[\chi]
  =
  \hat\chi_-^\dagger\hat\chi_+ ,
  \label{eq:doubled-gain-loss-generators}
\end{equation}
up to the fixed vectorization sign convention.  These generators conserve the
doubled charge \(N_++N_-\).  Since \(J_{\rm gain}^2=J_{\rm loss}^2=0\) on a
fixed mode, their finite-regulator Gaussian families are unipotent rather than
Pin reflections:
\begin{align}
  e^{\lambda J_{\rm gain}[\chi]}
  &:\quad
  \mathbb G_\lambda(\chi)
  =
  \I+\lambda\,\ket{-}\!\bra{+}\otimes P_\chi ,
  \label{eq:doubled-gain-unipotent}\\
  e^{\mu J_{\rm loss}[\chi]}
  &:\quad
  \mathbb L_\mu(\chi)
  =
  \I+\mu\,\ket{+}\!\bra{-}\otimes P_\chi ,
  \label{eq:doubled-loss-unipotent}
\end{align}
where \(P_\chi=\chi\chi^\dagger\).  At every finite \(\lambda,\mu\) these
compose by literal ordered matrix multiplication in the doubled one-particle
space.  The exact jump is the projective endpoint
\begin{equation}
  J_{\rm gain}[\chi]
  =
  \lim_{\lambda\to\infty}\lambda^{-1}e^{\lambda J_{\rm gain}[\chi]},
  \qquad
  J_{\rm loss}[\chi]
  =
  \lim_{\mu\to\infty}\mu^{-1}e^{\mu J_{\rm loss}[\chi]} .
  \label{eq:doubled-projective-jump-limit}
\end{equation}
This gives a second chart for the same exact endpoint and reproduces the
composition rule obtained by explicit Choi contraction.  It is not the same
finite matrix family as
\(\hat\chi^\dagger+\varepsilon\hat\chi\) or
\(\hat\chi+\varepsilon\hat\chi^\dagger\).

For full-rank even Choi states with particle transfer matrices \(T\) and
\(S\), contracting the right leg of the first with the left leg of the second
gives
\begin{equation}
  R=TS .
  \label{eq:choi-contraction-product}
\end{equation}
The same statement holds on the doubled \(+/-\) space with
\(T,S\) replaced by \(\mathbb T,\mathbb S\).  For an exact gain, the isolated
particle chart becomes rank deficient, \(T\mapsto TQ\) with
\(Q=\I-P_\chi\), but a subsequent Choi contraction with an even piece \(S\)
does not simply keep the undressed projector \(P_\chi\).  The contraction
derivation gives the dressed finite particle transfer
\begin{equation}
  R_{\rm gain}
  =
  TS\left(\I-P_{S^\dagger\chi}\right),
  \qquad
  P_{S^\dagger\chi}
  =
  \frac{S^\dagger P_\chi S}
       {\chi^\dagger SS^\dagger\chi}.
  \label{eq:choi-gain-dressed-projector}
\end{equation}
Thus the gain constraint is transported contravariantly through the later
factor \(S\).  The corresponding loss formula is most naturally read in the
hole chart:
\begin{equation}
  R_{{\rm loss},h}
  =
  (TS)^{-T}\left(\I-P_{S^{-1}\chi}^{\,T}\right),
  \qquad
  P_{S^{-1}\chi}
  =
  \frac{S^{-1}P_\chi S^{-\dagger}}
       {\chi^\dagger S^{-\dagger}S^{-1}\chi}.
  \label{eq:choi-loss-dressed-projector}
\end{equation}
Writing only the homogeneous particle-side expression for loss misses the
source term carried by the emptied endpoint sector.  This is precisely the
dictionary obstruction: gain is reconstructed cleanly in the particle chart,
loss in the hole chart, while mixed gain/loss composites can obstruct both
fixed charts and must be handled in a graph/covariance chart.  The doubled
transfer chart below is one stable implementation of that contraction, not an
additional assumption needed to derive the dressed projector.

The same-mode mixed product illustrates the relation to projective
measurements.  In the doubled chart,
\begin{equation}
  \mathbb G_\lambda(\chi)\mathbb L_\mu(\chi)
  =
  \begin{pmatrix}
    \I & \mu P_\chi\\
    \lambda P_\chi & \I+\lambda\mu P_\chi
  \end{pmatrix}_{+/-}.
  \label{eq:doubled-gain-loss-product}
\end{equation}
Gauss decomposition of this block matrix shows that, as
\(\lambda,\mu\to\infty\), the unipotent factors trivialize and the Levi block
degenerates to \(Q=\I-P_\chi\) on the appropriate copy.  This is the
Choi-transfer version of
\[
  \hat\chi\,\hat\chi^\dagger=\I-\hat n_\chi,
  \qquad
  \hat\chi^\dagger\hat\chi=\hat n_\chi .
\]
Therefore the apparent discrepancy is only a discrepancy if the two
regularizations are identified.  The single-copy Nambu regulator already
gives the dressed projector after the graph/covariance limit.  The
neutral/doubled Choi regulator is kept here only as a comparison chart in
which channel contraction is also literal matrix multiplication.

For cumulative operator spectroscopy, the clean representation is the Choi
state of the realized branch.  Let \(\ket{V}\!\rangle\) be the Choi vector of
an even Gaussian operator \(V\), and let \(\chi_R\) denote the right-leg mode.
The odd gain/loss amplitudes are represented by
\begin{equation}
  \ket{K_+}\!\rangle=\hat\chi_R^\dagger\ket{V}\!\rangle,
  \qquad
  \ket{K_-}\!\rangle=\hat\chi_R\ket{V}\!\rangle .
  \label{eq:odd-choi-branches}
\end{equation}
These are still pure Gaussian Choi states when their amplitudes are nonzero.
The complication is therefore not a failure of Gaussianity or purity; it is a
failure of the full-rank even particle-transfer chart.

Equivalently, exact right-leg gain/loss can be treated directly as occupation
constraints on the Choi covariance.  Let \(\mathcal C\) be the particle
correlation matrix of a normalized Choi Gaussian state on the doubled
left/right one-particle space, and let \(X_R\) be a matrix whose columns are
right-leg wavefunctions to be constrained, embedded in the doubled space with
zero support on the left leg.  The columns need not be orthonormal.  Put
\begin{equation}
  S_X=X_R^\dagger X_R,
  \qquad
  P_X=X_R S_X^{-1}X_R^\dagger
  \label{eq:choi-nonorthogonal-projector}
\end{equation}
when \(X_R\) has full column rank; if not, first remove the linearly dependent
columns.  A simultaneous gain constraint means that every mode in
\(\operatorname{span}X_R\) is filled.  Its raw ordered-operator norm is
\begin{equation}
  p^{\rm raw}_{+X}
  =
  \det\!\big[X_R^\dagger(\I-\mathcal C)X_R\big] .
  \label{eq:choi-gain-raw-gram}
\end{equation}
If instead one wants the basis-independent occupation constraint associated
only with the subspace \(P_X\), the normalized probability is
\begin{equation}
  p_{+X}
  =
  \frac{\det\!\big[X_R^\dagger(\I-\mathcal C)X_R\big]}{\det S_X} .
  \label{eq:choi-gain-normalized-gram}
\end{equation}
The conditional Choi covariance is independent of the chosen basis for the
constraint subspace:
\begin{equation}
  \mathcal C_{+X}
  =
  \mathcal C
  +
  (\I-\mathcal C)X_R
  \big[X_R^\dagger(\I-\mathcal C)X_R\big]^{-1}
  X_R^\dagger(\I-\mathcal C) .
  \label{eq:choi-gain-constraint-update}
\end{equation}
It satisfies
\begin{equation}
  \mathcal C_{+X}P_X=P_X,
  \qquad
  \Sigma_{+X}P_X=+P_X,
  \qquad
  \Sigma_{LR,+X}P_X=0,
  \label{eq:choi-gain-endpoint-constraint}
\end{equation}
where \(\Sigma=2\mathcal C-\I\).  Thus the right-leg constrained subspace is
an occupied endpoint sector.

The loss constraint is the particle-hole dual.  The raw and normalized weights
are
\begin{align}
  p^{\rm raw}_{-X}
  &=
  \det\!\big[X_R^\dagger\mathcal C X_R\big],
  \label{eq:choi-loss-raw-gram}\\
  p_{-X}
  &=
  \frac{\det\!\big[X_R^\dagger\mathcal C X_R\big]}{\det S_X},
  \label{eq:choi-loss-normalized-gram}
\end{align}
and, when the occupied Gram matrix is nonsingular,
\begin{equation}
  \mathcal C_{-X}
  =
  \mathcal C
  -
  \mathcal C X_R
  \big[X_R^\dagger\mathcal C X_R\big]^{-1}
  X_R^\dagger\mathcal C .
  \label{eq:choi-loss-constraint-update}
\end{equation}
Now
\begin{equation}
  \mathcal C_{-X}P_X=0,
  \qquad
  \Sigma_{-X}P_X=-P_X,
  \qquad
  \Sigma_{LR,-X}P_X=0,
  \label{eq:choi-loss-endpoint-constraint}
\end{equation}
so the right-leg constrained subspace is an empty endpoint sector.  Equations
\eqref{eq:choi-gain-constraint-update} and
\eqref{eq:choi-loss-constraint-update} are often the cleanest exact way to
handle gain and loss in numerics: impose a right-leg occupation constraint on
the Choi covariance, then read the finite transfer block only on the
orthogonal complement of \(P_X\).

Non-orthogonal wavefunctions cause no conceptual problem for same-type
constraints.  They only replace rank-one denominators by Gram determinants and
Schur complements.  What matters is the rank and conditioning of
\(X_R^\dagger(\I-\mathcal C)X_R\) for gains or
\(X_R^\dagger\mathcal C X_R\) for losses.  A zero determinant means a blocked
or redundant exact branch.  It should be reported as endpoint/null-sector
bookkeeping, not used as a finite transfer gap.

Mixed gain and loss constraints are different.  A filled subspace and an empty
subspace can be imposed simultaneously only when they are compatible; in
particular, overlapping non-orthogonal directions cannot both be definitely
filled and definitely empty.  For example
\(\hat\beta\hat\alpha^\dagger=\langle\beta|\alpha\rangle-
\hat\alpha^\dagger\hat\beta\), so changing the order changes the scalar
contraction and the residual operator.  The correct Choi-covariance procedure
for a mixed non-commuting word is therefore the ordered Pfaffian/Schur
conditioning of Sec.~\ref{sec:several-odd-insertions}, applied on the doubled
Choi state with all operators embedded on the appropriate leg.  Only after
that ordered conditioning has been performed does one identify endpoint
multiplicity and the remaining finite Choi transfer spectrum.

To see this in the smallest possible example, take a single mode and \(V=\I\).
Both right-leg creation and right-leg loss have vanishing effective bridge
between left and right particle sectors,
\begin{equation}
  T_{\rm eff}=0 .
  \label{eq:single-mode-Teff-zero}
\end{equation}
Nevertheless they produce different Choi states:
\begin{equation}
  \hat\chi_R^\dagger\ket{\Omega}\propto\ket{1_L1_R},
  \qquad
  \hat\chi_R\ket{\Omega}\propto\ket{0_L0_R}.
  \label{eq:single-mode-odd-completions}
\end{equation}
The singular matrix \(T_{\rm eff}=0\) alone does not know whether the
disconnected null sector is filled or empty.

For an isolated right-leg insertion after an even operator with particle
transfer \(T_p\), the rank-deficient particle-chart propagation part is
\begin{equation}
  T_p \longmapsto T_p(\I-P).
  \label{eq:right-gain-loss-column-removal}
\end{equation}
This is not the full Choi-contraction composition rule once later operators
are present.  If the constrained Choi state is contracted with a subsequent
even transfer \(S\), the projector is dressed as in
Eqs.~\eqref{eq:choi-gain-dressed-projector} and
\eqref{eq:choi-loss-dressed-projector}; gain transports
\(\chi\) by \(S^\dagger\), while loss transports it by \(S^{-1}\) and should
be reconstructed in the hole chart.
The complementary hole chart obeys \(T_h=(T_p^\dagger)^{-1}\) before the
singular limit.  Hence an exact particle-dark direction is a pole in the hole
description, and an exact hole-dark direction is a pole in the particle
description.  The finite-sector spectrum must therefore be supplemented by
endpoint data:
\begin{equation}
  \left\{\lambda_{p,j}^{\rm finite}\right\}
  \quad {\rm plus}\quad
  \left(P_{\rm null},n_{\rm null},m_{\rm null}\right),
  \label{eq:complete-singular-data}
\end{equation}
where \(P_{\rm null}\) is the projected direction, \(n_{\rm null}=1\) for
creation and \(n_{\rm null}=0\) for annihilation, and \(m_{\rm null}\) is the
endpoint multiplicity.

In the Choi-covariance notation \(\Sigma=2C_{\rm Choi}-\I\), the finite
particle exponents can be read from the left-left block eigenvalues
\(a_j\in(-1,1)\):
\begin{equation}
  \lambda_{p,j}(t)
  =
  \frac{1}{2t}
  \log\!\left(\frac{1-a_j(t)}{1+a_j(t)}\right)
  =
  -\frac{1}{t}\artanh a_j(t),
  \qquad
  a_j(t)\in\spec\Sigma_{LL}(t)\cap(-1,1).
  \label{eq:particle-exponents-from-choi}
\end{equation}
Endpoint sectors \(a_j=\pm1\) are formal infinite exponents in one of the
particle/hole charts.  A finite Lyapunov gap should therefore be defined as
\begin{equation}
  \Delta_\lambda^{\rm finite}(t)
  =
  \min_{a_j(t)\in(-1,1)}
  \left|
  \frac{1}{t}\artanh a_j(t)
  \right|,
  \label{eq:finite-sector-gap}
\end{equation}
with the number of excluded endpoint eigenvalues reported separately.  This is
the gain/loss bookkeeping required before comparing transfer gaps across
protocols or sizes.


\subsection{Operational Conclusion for Non-Commuting Gain/Loss Strings}
\label{sec:gain-loss-operational-conclusion}

The practical rule is the following.  If the trajectory contains a block of
gains and losses, first read the word in chronological order and form the
running charge \(Q_j\) of Eq.~\eqref{eq:running-charge-definition}.  Cut at
returns to the running minimum.  Each neutral block with nonzero causal Gram
determinant gives an honest finite graph transfer through
Eq.~\eqref{eq:balanced-block-master-formula}; its scalar
\(\lambda=\pm\det\widetilde G\) is recorded in the Born-weight cocycle.  Each
unbalanced block is endpoint orbitals plus one finite matrix, as in
Eq.~\eqref{eq:unbalanced-gain-block-normal-form}, not a matrix by itself.

For same-type blocks this reduces to ordinary occupation constraints.  Gains
use the empty Gram matrix \(X^\dagger(\I-C)X\), update the covariance by
Eq.~\eqref{eq:many-gain-covariance}, and record the rank of the filled
constrained subspace as endpoint data.  Losses use the occupied Gram matrix
\(Y^\dagger C Y\), update by Eq.~\eqref{eq:many-loss-covariance}, and record
the rank of the emptied constrained subspace.  For mixed non-commuting
blocks, do not first combine all gains and all losses.  Either apply the
one-site gain/loss updates sequentially in chronological order, or build the
ordered Pfaffian problem
Eqs.~\eqref{eq:mixed-kdagger-k-list}--\eqref{eq:mixed-pfaffian-schur-update};
these are the same Wick calculation written in different forms.

In a fixed particle or hole transfer chart, the finite transfer matrix is
therefore a reduced object obtained after the ordered constraints have been
resolved.  Its endpoint multiplicity is a rank computed from the ordered
constrained subspace.  It is not, in general, the number of gain and loss
events.  Mixed gain/loss pairs can annihilate into projectors, produce scalar
contractions through Eq.~\eqref{eq:exterior-interior-clifford-relations}, or
leave genuine charge-changing exterior/interior maps.  The many-body spectrum
must be built from the resulting ordered Fock-space map.

In a moving Majorana chart, the equivalent prescription is to replace each odd
insertion by its regularized Pin representative, multiply all associated
\(O(2N,\mathbb C)\) matrices in chronological order, and accumulate the
projective scalar normalization separately.  The singular limit
\(\varepsilon\to0\) must be taken only after the ordered product and scalar
weight have been formed.  The matrix product before the limit is literally
ordinary matrix multiplication; the limiting finite object is the
graph/covariance block in Eq.~\eqref{eq:graph-block-convergence}, with
endpoint/null sectors retained separately.  This preserves the correct Born
probability and prevents the divergent gauge directions of the regularized
chart from being mistaken for finite CFT Lyapunov gaps.

\section{Tangent Cocycles for Covariance Dynamics}
\label{sec:tangent-cocycle}

The state covariance evolves by a branch-dependent nonlinear map
\begin{equation}
  G_{t+1}=F_{\alpha_t}(G_t),
  \qquad
  \alpha_t\sim p_t(\cdot\mid G_t).
  \label{eq:branch-markov-map}
\end{equation}
For a Schur-complement Gaussian branch, a convenient form is
\begin{equation}
  F_\alpha(G)
  =
  C_\alpha
  -
  B_\alpha^\dagger
  \left(A_\alpha-G^{-1}\right)^{-1}
  B_\alpha .
  \label{eq:fractional-branch-map}
\end{equation}
Differentiating gives the tangent map.  Let
\begin{equation}
  M_\alpha(G)=A_\alpha-G^{-1}.
  \label{eq:M-alpha}
\end{equation}
Since
\begin{equation}
  \delta G^{-1}=-G^{-1}\delta G\,G^{-1},
  \qquad
  \delta M_\alpha=G^{-1}\delta G\,G^{-1},
  \label{eq:inverse-differentiation}
\end{equation}
and \(\delta M^{-1}=-M^{-1}(\delta M)M^{-1}\), one obtains
\begin{align}
  D F_\alpha(G)[\delta G]
  &=
  B_\alpha^\dagger
  M_\alpha(G)^{-1}
  G^{-1}\delta G\,G^{-1}
  M_\alpha(G)^{-1}
  B_\alpha .
  \label{eq:DF-first-form}
\end{align}
Using
\begin{equation}
  M_\alpha(G)^{-1}G^{-1}=(G A_\alpha-\I)^{-1},
  \qquad
  G^{-1}M_\alpha(G)^{-1}=(A_\alpha G-\I)^{-1},
  \label{eq:noncommuting-identities}
\end{equation}
define
\begin{equation}
  T_\alpha(G)=(A_\alpha G-\I)^{-1}B_\alpha .
  \label{eq:tangent-T}
\end{equation}
When the Hermiticity relations of the covariance contraction identify the left
factor with the adjoint of the right factor,
\begin{equation}
  \delta G_{t+1}
  =
  T_{\alpha_t}(G_t)^\dagger\delta G_tT_{\alpha_t}(G_t).
  \label{eq:tangent-update}
\end{equation}

Along a realized trajectory set \(T_t=T_{\alpha_t}(G_{t-1})\).  The
chronological product convention used by Eq.~\eqref{eq:tangent-update} is
\begin{equation}
  \widetilde T_n=T_1T_2\cdots T_n,
  \qquad
  \delta G_n=\widetilde T_n^\dagger\delta G_0\widetilde T_n .
  \label{eq:tangent-product}
\end{equation}
The single-particle tangent Lyapunov exponents are
\begin{equation}
  \lambda_i^{\rm tan}
  =
  \lim_{n\to\infty}
  \frac{1}{n}\log\sigma_i(\widetilde T_n),
  \label{eq:tangent-lyapunov}
\end{equation}
assuming the stationary ergodic hypotheses of the multiplicative ergodic
theorem~\cite{Oseledets1968,FurstenbergKesten1960}.  These are quenched
singular-value growth rates along sampled trajectories.  In general,
\begin{equation}
  \E\!\left[T_\alpha(G)^\dagger\delta G\,T_\alpha(G)\mid G\right]
  \neq
  \bar T(G)^\dagger\delta G\,\bar T(G),
  \qquad
  \bar T(G)=\E[T_\alpha(G)\mid G].
  \label{eq:no-average-transfer}
\end{equation}
Thus the tangent Lyapunov spectrum is not the spectrum of an outcome-averaged
transfer matrix.

\section{Trajectory Free Energy and CFT Finite-Size Data}
\label{sec:cft-extraction}

For a monitored circuit on a cylinder of circumference \(L\) and time length
\(t\), define the Born-weighted trajectory free energy
\begin{equation}
  F(L,t)
  =
  -\sum_\bxi p_\bxi\log p_\bxi
  =
  -\E_{\bxi\sim p}\log p_\bxi .
  \label{eq:trajectory-free-energy}
\end{equation}
In a Gaussian circuit, each \(p_\bxi\) is computed by multiplying the
determinant branch normalizations derived above.  If the intrinsic
space-time anisotropy is \(\alpha\), define
\begin{equation}
  f_0(L)=\lim_{t\gg L}\frac{F(L,t)}{\alpha Lt}.
  \label{eq:free-energy-density}
\end{equation}
At a conformally invariant non-unitary fixed point, the cylinder free-energy
correction has the universal form~\cite{Cardy1986,DiFrancesco1997}
\begin{equation}
  f_0(L)=f_0(\infty)-\frac{\pi c_{\rm eff}}{6L^2}+\cdots .
  \label{eq:ceff-cylinder}
\end{equation}
Here \(c_{\rm eff}\) is the replica effective central charge relevant to the
trajectory ensemble, not automatically the entanglement central charge of a
prepared edge state.

\subsection{Why Transfer Eigenvalues Are Free Energies}
\label{sec:transfer-free-energy-identity}

The transfer-matrix statement is just the row decomposition of a two-dimensional
statistical mechanics sum.  For a fixed circumference \(L\), collect all local
Boltzmann weights, Born weights, or trajectory amplitudes in one time strip
into a linear operator \(\cT_L\) acting on the Hilbert space of boundary data
around the cylinder.  Gluing two neighboring strips means summing over their
shared boundary variables, which is matrix multiplication.  Gluing \(t\)
strips and finally closing the time direction gives
\begin{equation}
  Z(L,t)
  =
  \Tr_{\rm bdry}\,\cT_L^t .
  \label{eq:row-transfer-partition-function}
\end{equation}
If \(\cT_L\) is diagonalizable with eigenvalues \(\Lambda_A(L)\), this is
\begin{equation}
  Z(L,t)=\sum_A \Lambda_A(L)^t .
  \label{eq:transfer-eigenvalue-sum}
\end{equation}
For a non-Hermitian or trajectory-random problem one often computes singular
values or Lyapunov exponents of products instead of ordinary eigenvalues, but
the meaning is the same: the relevant number is the logarithmic growth or
decay rate per time strip.  Projecting onto, or inserting an operator that
selects, a sector \(A\) gives
\begin{equation}
  Z_A(L,t)\sim \Lambda_A(L)^t,
  \qquad
  F_A(L,t)\equiv-\log Z_A(L,t)
  =
  -t\log \Lambda_A(L)+O(1),
  \label{eq:sector-transfer-growth}
\end{equation}
where \(\log\Lambda_A\) should be read as \(\log|\Lambda_A|\) when phases do
not affect the decay rate.  Dividing by the cylinder area \(\alpha Lt\)
therefore gives
\begin{equation}
  f_A(L)
  =
  -\frac{\log|\Lambda_A(L)|}{\alpha L},
  \qquad
  f_A(L)-f_0(L)
  =
  \frac{-\log|\Lambda_A(L)|+\log|\Lambda_0(L)|}{\alpha L}.
  \label{eq:transfer-eigenvalue-free-energy-gap}
\end{equation}
This is the precise sense in which subleading transfer eigenvalues, singular
values, or quenched Lyapunov exponents are ``free-energy gaps.''  The
Lyapunov calculation is an efficient numerical route to the logarithms in
Eq.~\eqref{eq:transfer-eigenvalue-free-energy-gap}; it is not an additional
CFT postulate.

\subsection{Replicas Are Not a Numerical Requirement}
\label{sec:replica-clarification}

The numerical estimators in this note do not require introducing explicit
replica degrees of freedom.  A sampled trajectory directly supplies
\(-\log p_\bxi\), and the typical gap is obtained by taking a logarithm of the
trajectory-resolved transfer growth before averaging over \(\bxi\).  Replicas
are instead an analytic device: they rewrite quenched averages and trajectory
moments as limits or continuations of integer-replica statistical mechanics
models.  The notation \(c_{\rm eff}\) refers to the effective central charge
of that non-unitary trajectory ensemble.  It should not be identified with the
entanglement central charge \(c=1\) of a prepared chiral free-fermion edge
unless a separate argument relates those two conformal theories.

Scaling dimensions are extracted from many-body transfer gaps, not directly
from an arbitrary one-body Lyapunov spectrum.  Let
\(\widehat{\mathcal T}_\bxi(t)\) denote the unnormalized trajectory transfer
operator in the sector whose correlation function is being probed, and let
\(S_A[\widehat{\mathcal T}_\bxi(t)]\) be its many-body singular values ordered
by decreasing typical growth rate.  Define the quenched many-body gaps by
\begin{equation}
  \mu_A(L)
  =
  -\lim_{t\to\infty}
  \frac{1}{t}
  \E_\bxi
  \log
  \frac{
    S_A[\widehat{\mathcal T}_\bxi(t)]
  }{
    S_0[\widehat{\mathcal T}_\bxi(t)]
  },
  \qquad
  \mu_A(L)\ge 0 .
  \label{eq:many-body-transfer-gap}
\end{equation}
The associated generalized free-energy gap is
\begin{equation}
  \Delta f_A(L)
  =
  \frac{\mu_A(L)}{\alpha L}.
  \label{eq:many-body-free-energy-gap}
\end{equation}
For low-lying operators of the effective CFT,
\begin{equation}
  \Delta f_A(L)
  =
  \frac{2\pi x_A^{\rm typ}}{L^2}+\cdots ,
  \qquad
  \mu_A(L)
  =
  \frac{2\pi\alpha x_A^{\rm typ}}{L}+\cdots .
  \label{eq:scaling-dimensions-from-gaps}
\end{equation}
The superscript ``typ'' emphasizes that the logarithm is taken before the
trajectory average.

For free fermions, Eq.~\eqref{eq:many-body-transfer-gap} can often be reduced
to one-body data, but only after reconstructing the exterior-power or
Fock-space tower of Sec.~\ref{sec:one-body-many-body-spectra}.  If the
realized even Gaussian transfer object has finite one-body excitation
rapidities \(\epsilon_a(L)\), with endpoint/null sectors removed as in
Sec.~\ref{sec:gain-loss}, then many-body gaps in a fixed charge/parity sector
are additive combinations
\begin{equation}
  \mu_I(L)
  =
  \sum_a n_a(I)\,\epsilon_a(L),
  \qquad
  n_a(I)\in\mathbb Z_{\ge 0},
  \label{eq:fock-tower-rapidities}
\end{equation}
up to the scalar normalization common to all states in the sector.  The
smallest positive one-body rapidity equals the first many-body gap only when
the one-body list has already been measured relative to the dominant
many-body configuration and the corresponding single-quasiparticle excitation
is allowed by the sector and by the observable being measured.  Otherwise the
first CFT gap may be a difference of many-body sums, a sum of several
one-body excitation rapidities, or a gap in a different particle/hole chart.

The tangent cocycle of Sec.~\ref{sec:tangent-cocycle} is also first obtained
as a one-body cocycle \(T_\bxi\).  Its associated Gaussian many-body tangent
transfer spectrum is again the additive Fock-space lift of the one-body
Lyapunov data.  Thus the raw \(\lambda^{\rm tan}_a\) are not the final
many-body gaps; the gaps are sums
\begin{equation}
  \mu_I^{\rm tan}(L)
  =
  \sum_a n_a(I)\,\epsilon^{\rm tan}_a(L),
  \qquad
  \epsilon^{\rm tan}_a(L)\equiv
  \bigl|\lambda^{\rm tan}_a(L)-\lambda^{\rm tan}_0(L)\bigr|_{\rm finite},
  \label{eq:tangent-fock-tower}
\end{equation}
with the sector \(I\) fixed by the perturbation or correlator being tested.
Vectorizing \(\delta G\) is a useful way to implement the linear covariance
map, but it does not by itself choose the many-body exterior-power sector or
the excitation pattern \(n_a(I)\).

Average correlations need not share the same exponent.  If a trajectory
correlator satisfies
\begin{equation}
  G_\bxi(t)\sim
  \exp\!\left[-\frac{2\pi t}{L}x_{\rm typ}\right]
  \label{eq:typical-correlator}
\end{equation}
typically, its moments may instead scale as
\begin{equation}
  \E\!\left[G_\bxi(t)^n\right]
  \sim
  \exp\!\left[-\frac{2\pi t}{L}x(n)\right].
  \label{eq:moment-exponents}
\end{equation}
Writing \(Y=-\log G_\bxi(t)\), multifractal scaling is the statement that
the distribution of \(Y\) has the large-deviation form~\cite{EversMirlin2008}
\begin{equation}
  P_t(Y)
  \sim
  \left(\frac{2\pi\alpha t}{L}\right)^{-1/2}
  \exp\!\left[
    -\frac{2\pi\alpha t}{L}
    H\!\left(\frac{Y}{2\pi\alpha t/L}\right)
  \right],
  \label{eq:multifractal-distribution}
\end{equation}
with a nontrivial scaling function \(H\).  Expanding cumulants gives, for
small \(n\),
\begin{equation}
  x(n)=n x^{(1)}+\frac{n^2}{2}x^{(2)}+\cdots ,
  \qquad
  x^{(1)}=x_{\rm typ}.
  \label{eq:cumulant-expansion}
\end{equation}
The operational lesson is that a Lyapunov gap, an average correlator, and a
typical correlator may be three different finite-size measurements of related
but inequivalent universal data.

\section{Mechanical CPU Extraction Plan}
\label{sec:practical-protocol}

The numerical extraction must be written so that every accumulated number has
a named transfer object.  For the domain-wall runs below the cylinder
circumference is
\begin{equation}
  L \equiv N_y ,
  \label{eq:cpu-circumference}
\end{equation}
and the anisotropy is reported with the default convention \(\alpha=1\) unless
an independent calibration is supplied.  The production CPU driver samples the
physical Born circuit only:
\begin{equation}
  {\tt postselect}={\tt False},
  \qquad
  {\tt postselect\_probability}=0 .
  \label{eq:cpu-no-postselect}
\end{equation}
Forced postselection is useful as a diagnostic trajectory, but it does not
define the Born ensemble weight \(p_\bxi\) and is therefore excluded from the
\(c_{\rm eff}\) estimator.

For sample \(s\), size \(L\), and completed cycle \(t\), the CPU observer
records the site-branch probabilities \(p_{s,\ell}\) generated by the
canonical covariance update.  These include the Born probability of the
measured occupation outcome and, when a stochastic correction is used, the
probability of the sampled gain/loss correction.  Deterministic perfect
correction contributes unit probability.  The accumulated trajectory action is
\begin{equation}
  Y_s(L,t)
  =
  -\log p_{\bxi_s(0:t)}
  =
  -\sum_{\ell\in[0,t)}\log p_{s,\ell}.
  \label{eq:cpu-trajectory-action}
\end{equation}
The free-energy density estimator is then
\begin{equation}
  f_0(L,t)
  =
  \frac{1}{\alpha L t}
  \frac{1}{N_{\rm samp}(L)}
  \sum_{s=1}^{N_{\rm samp}(L)}Y_s(L,t).
  \label{eq:cpu-f0-estimator}
\end{equation}
After discarding early-time data and choosing a plateau time window, the
cylinder fit is
\begin{equation}
  f_0(L)
  =
  f_\infty
  -
  \frac{\pi c_{\rm eff}}{6L^2}
  +
  O(L^{-4}).
  \label{eq:cpu-ceff-fit}
\end{equation}
Thus a linear fit of \(f_0(L)\) against \(1/L^2\) with slope \(m_0\) gives
\begin{equation}
  c_{\rm eff}=-\frac{6m_0}{\pi}.
  \label{eq:cpu-ceff-from-slope}
\end{equation}

In the same physical trajectories the code accumulates two one-body Lyapunov
objects and then builds the corresponding additive many-body gaps.  The first
is the covariance tangent cocycle,
\begin{equation}
  \delta G_{t+1}
  =
  T_{\alpha_t}(G_t)^\dagger
  \delta G_t
  T_{\alpha_t}(G_t),
  \label{eq:cpu-tangent-object}
\end{equation}
stabilized by QR decomposition after each cycle.  The resulting exponents are
called \(\lambda^{\rm tan}_{j,s}(L,t)\).  The second is the Choi/operator
covariance transfer object with the endpoint bookkeeping of
Sec.~\ref{sec:gain-loss}; diagonalizing the finite particle/hole block after
removing exact gain/loss endpoints gives
\(\epsilon^{\rm Choi}_{j,s}(L,t)\).  These are one-body rapidities from which
many-body tangent and Choi/Kraus gaps must both be assembled additively, as in
Eqs.~\eqref{eq:tangent-fock-tower} and
\eqref{eq:fock-tower-rapidities}.  The analysis therefore reports
\begin{equation}
  \delta\ell_{\rm tan}^{(1)}
  \qquad
  \epsilon_{\rm Choi}^{(1)}
  \qquad
  x_{\rm typ}^{\rm tan/Fock}
  \qquad
  x_{\rm typ}^{\rm Choi/Fock},
  \label{eq:three-spectrum-labels}
\end{equation}
where the first two are raw one-body ingredients and the last two are CFT
scaling-dimension estimates only after choosing allowed many-body Fock
excitations.

For either one-body route \(r\in\{\mathrm{tan},\mathrm{Choi}\}\), the finite
single-particle gap recorded by the CPU code is
\begin{equation}
  \delta\ell^{r}_{k}(L,t)
  =
  {\rm median}_{s}
  \left[
    {\rm kth\,smallest}_{j}\,
    \bigl|\ell^{r}_{j,s}(L,t)\bigr|
  \right]_{\rm finite},
  \label{eq:cpu-one-body-gap}
\end{equation}
with \(\ell^{\rm tan}=\lambda^{\rm tan}\) and
\(\ell^{\rm Choi}=\epsilon^{\rm Choi}\).  Here the Choi rapidities are obtained
from the finite covariance eigenvalues by the particle/hole chart convention,
while the tangent exponents are obtained from QR-stabilized tangent-frame
growth.
The finite bracket means that zero/floor directions, endpoint directions,
nonfinite values, and censored Choi trajectories are excluded before ordering.
Their multiplicities are retained in the output table.

For each route the next mechanical step is to build the allowed many-body gaps
\begin{equation}
  \mu_A^{r}(L,t)
  =
  \sum_j n_j^{r}(A)\,\delta\ell^{r}_{j}(L,t),
  \qquad
  r\in\{\mathrm{tan},\mathrm{Choi}\},
  \label{eq:cpu-fock-gap}
\end{equation}
with occupation pattern \(A\) fixed by the operator sector being tested.  Only
these \(\mu_A^r\), not the raw one-body list by itself, enter
the CFT fit
\begin{equation}
  \frac{\mu_A^{r}(L,t)}{\alpha L}
  =
  \frac{2\pi x_A^{\rm typ}}{L^2}
  +
  O(L^{-3}).
  \label{eq:cpu-fock-x-fit}
\end{equation}
The often-used estimator
\begin{equation}
  x_{\rm one\ body}^{r}(L,t)
  =
  \frac{L\,\delta\ell^{r}_{1}(L,t)}{2\pi\alpha}
  \label{eq:cpu-one-body-estimator}
\end{equation}
is therefore a CFT scaling dimension only under the additional assumption that
the first allowed many-body excitation is exactly one finite quasiparticle in
route \(r\).
The CPU script parameter \({\tt fock\_rank}=R\) implements only the mechanical
choice
\begin{equation}
  n_j^r(A_R)=
  \begin{cases}
    1, & 1\le j\le R,\\
    0, & j>R,
  \end{cases}
  \qquad
  \mu_{A_R}^{r}(L,t)=\sum_{j=1}^{R}\delta\ell_j^r(L,t),
  \label{eq:cpu-fock-rank-convention}
\end{equation}
after ordering the finite one-body gaps.  This is a transparent low-lying
Fock-tower convention, not a claim that the physical operator of interest
necessarily occupies the first \(R\) finite quasiparticle modes.
If the sector requires a pair, a fixed charge difference, or a particle-hole
combination, Eq.~\eqref{eq:cpu-fock-gap} must be used instead.  Agreement
between the tangent/Fock and Choi/Fock dimensions is evidence that the same
slow channel is being observed by two Gaussian transfer constructions, not a
proof by notation.

Gain and loss introduce the most important mechanical caveat.  In the Choi
route, eigenvalues satisfying
\begin{equation}
  \bigl||g_j|-1\bigr|<\varepsilon_{\rm end}
  \label{eq:cpu-endpoint-test}
\end{equation}
are endpoint sectors and are not used as finite rapidity gaps.  In the tangent
route, QR directions whose diagonal factors hit the numerical floor, or whose
cycle exponents become nonfinite, are classified as null/divergent sectors.
Both endpoint counts are physical data: they measure how many rank-one
creation/removal resets have left the chosen particle or hole chart.  They are
reported beside the finite-sector fits, not thrown away silently.

The concrete CPU workflow is therefore: run the canonical
\({\tt classA\_U1FGTN.run\_markov\_circuit}\) entry point with serial samples,
enable the trajectory-weight observer, tangent observer, and Choi observer,
save per-size raw spectra and probability actions, bootstrap over samples,
fit \(f_0(L)\) against \(1/L^2\) for \(c_{\rm eff}\), plot the tangent
and Choi one-body spectra, and convert both to \(x_{\rm typ}\) only through
the specified additive many-body Fock excitation.

For the current domain-wall adaptive circuit, the conservative interpretation
is therefore: the observed small transfer-frame gap is evidence for a
near-marginal interface channel localized on the wall.  It is not, by itself,
a precision extraction of \(x_i\) or \(c_{\rm eff}\).  Such an extraction
requires a controlled sequence in \(L\), endpoint bookkeeping for gain/loss,
and a clear choice of operator sector.

\section{Caveats and Non-Identifications}
\label{sec:caveats}

\begin{enumerate}
  \item The Born-sampled trajectory ensemble is not the deterministic mean
  channel.  For nonlinear observables \(O\),
  \begin{equation}
    \E_\bxi[O(G_\bxi)]\neq O(\E_\bxi[G_\bxi]).
    \label{eq:nonlinear-average-warning}
  \end{equation}

  \item A quenched product of random transfer matrices is not an averaged
  transfer matrix.  Equation~\eqref{eq:no-average-transfer} is the minimal
  algebraic obstruction.

  \item The Choi/Kraus particle spectrum and the tangent-cocycle spectrum are
  different linearizations.  The former can generate many-body Gaussian
  transfer gaps after Fock-space reconstruction; the latter is a covariance
  stability spectrum unless tied to a specific trajectory correlator.  Agreement
  between their slow sectors is physical evidence, not an operator identity.

  \item Gain and loss are not harmless details.  Exact creation or annihilation
  branches move the transfer description onto a singular boundary where a
  finite matrix must be supplemented by null-sector occupation data.

  \item The entanglement central charge \(c=1\) of a chiral free-fermion edge is
  not automatically the same as the non-unitary effective central charge
  \(c_{\rm eff}\) obtained from trajectory free energies.  The former is a
  property of the prepared state; the latter is a property of the trajectory
  ensemble viewed as a replicated statistical mechanics problem.
\end{enumerate}

\section{Summary}

In a free-fermion monitored circuit, the trajectory partition function
\(Z_\bxi=p_\bxi\) becomes a determinant calculation, and the many-body
trajectory singular spectrum can be reconstructed from additive one-body
singular data whenever the branch remains in a regular even Gaussian chart.
Measurements, gain, and loss keep the state Gaussian but introduce singular
transfer directions.  The correct spectroscopy is therefore not simply
``diagonalize a matrix'': one must state which matrix, which particle or hole
chart, which endpoint sectors have been removed, which Fock-sector excitation
is allowed, and whether the result concerns a cumulative operator or the
tangent dynamics of covariance perturbations.  Once that bookkeeping is
explicit, the finite-size CFT formulas for \(c_{\rm eff}\), scaling dimensions,
and multifractal spectra can be applied as a controlled diagnostic rather than
as a slogan.

\bibliographystyle{apsrev4-2}
\bibliography{}

\end{document}
